% !TeX root = main.tex
\chapter{Методи та інформаційна технологія виявлення ознак аномального вмісту в мультимедійних даних вебсередовища}
\label{chap:methods_technology}

Модель, розроблена в розділі~\ref{chap:models_formalization}, визначає склад представлень мультимедійного об’єкта та аналітичного стану, однак не встановлює способу їх одержання у процесі перевірки. Методичне завдання полягає в тому, щоб сформувати структурні ознаки з незмінної байтової послідовності, спектральні й піксельні ознаки --- з контрольовано розкодованих даних, а ознаки подій і поведінки --- з перебігу оброблення, не втративши спільного походження цих відомостей. Без такого зв’язку результати характеризуватимуть різні стани об’єкта і не утворюватимуть належної підстави для одного рішення.

Для розв’язання зазначеного завдання модель конкретизовано двома методами та інформаційною технологією. Метод аналізу спектральних, піксельних і структурних ознак формує статичний опис без передчасного змішування фізично різних ознак. Метод поведінкового узгодження ознак різного походження встановлює належність подій обробленню того самого об’єкта, після чого поєднує їх зі статичним описом. Інформаційна технологія визначає порядок ізоляції, аналізу, прийняття рішення, реагування та реєстрації.

Наукова відмінність одержаних результатів полягає не в повторному застосуванні відомих дискретного косинусного й дискретного багатомасштабного перетворень, аналізу молодших бітів, ентропійних показників, перевірок контейнера, маскування недоступних модальностей або послідовного статичного й динамічного аналізу. Кожний із цих засобів має відомі аналоги. Відмінність роботи обмежено доменною організацією цих засобів: збереженням спільного походження байтових, піксельних і подієвих свідчень, явними координатами й станами доступності, а також відокремленням аналітичного стану від технологічної дії. Межі відмінності від конкретних найближчих аналогів і стан перевірки кожного результату систематизовано в табл.~\ref{tab:ch3_novelty_comparison}. Кількісні результати наведено лише для складників, перевірених у розділі~\ref{chap:experiments}.

\begin{table}[p]
\centering
\footnotesize
\caption{Відмінність запропонованих результатів від найближчих класів аналогів}
\label{tab:ch3_novelty_comparison}
\begin{tabularx}{\textwidth}{@{}>{\raggedright\arraybackslash}p{0.15\textwidth}>{\raggedright\arraybackslash}p{0.20\textwidth}Y>{\raggedright\arraybackslash}p{0.20\textwidth}@{}}
\toprule
Результат & Найближчий клас аналогів & Відмінний елемент & Стан перевірки \\
\midrule
Модель мультимодального представлення & Карти походження мультимедійних об’єктів~\cite{kiltz2024trace,vielhauer2025modeling} і методи роботи з пропущеними модальностями~\cite{ma2021smil} & Формально визначене представлення, що в одному циклі пов’язує незмінні байти, похідні представлення, координати, причинно відібрані події та аналітичний стан & Сформульовано умови формальної узгодженості; повну сукупність складників не перевірено \\
Метод аналізу спектральних, піксельних і структурних ознак & Стегоаналітичні засоби та системи, що поєднують структурні й статистичні перевірки~\cite{badar2025stegomalware,kadebu2023hybrid}, а також аналіз контейнера~\cite{koch2022polyglot} & Роздільне формування двох статичних складових зі збереженням байтових і просторово-часових координат, доступності й форматного правила до утворення спільного опису & Детерміновану статичну конфігурацію перевірено на внутрішніх і зовнішніх наборах; ізольовану складову спектрального та піксельного аналізу зовні не підтверджено \\
Метод поведінкового узгодження ознак різного походження & Послідовний статичний і динамічний аналіз stegomalware~\cite{kadebu2023hybrid,badar2025stegomalware} та моделювання його комунікаційних трас~\cite{vielhauer2025modeling} & Причинно-фазове правило допуску події для того самого об’єкта й запуску, посилання на початковий запис і розрізнення порожнього та недоступного журналу & Формально визначено; причинно пов’язаних журналів і перевірки повної конфігурації немає \\
Інформаційна технологія, що набула подальшого розвитку & Ізольовані послідовності статичного й динамічного аналізу~\cite{kadebu2023hybrid,abbadini2023sandboxing} і реконструкція мультимедійного вмісту~\cite{belkind2023icdr} & Переходи між станами з формально визначеним складом, виконання зовнішнього припису лише після реєстрації та подання реконструкції як нового циклу з власним ідентифікатором & Локально перевірено переходи, атомарну реєстрацію, відмови, повтор і відновлення; реконструкцію та промислову інтеграцію не перевірено \\
\bottomrule
\end{tabularx}
\end{table}

Матеріал розділу викладено відповідно до переходу від загальної організації процесу до конкретних способів обчислення. Спочатку встановлено умови передавання результатів між етапами. Далі описано формування статичного представлення, його структурної складової та складової спектральних і піксельних ознак, після чого визначено порядок причинного допуску подій і їх узгодження зі статичними ознаками. Завершальний підрозділ установлює правила реагування, реєстрації та повторної перевірки реконструйованого об’єкта.

\section{Загальна організація процесу виявлення ознак аномального вмісту}
\label{sec:process_organization}

Загальний процес має розв’язати суперечність між необхідністю виконувати перетворення та вимогою зберегти початковий стан носія. Розкодування потрібне для одержання піксельних і часових ознак, проте воно може відкинути невідомі ділянки, унормувати службові значення або змінити порядок доступних даних. Реконструкція навмисно усуває частину початкової структури, тому не може передувати її перевірці. Відповідно, початковий мультимедійний об’єкт $X_{mm}$ зберігають незмінним, а всі операції, здатні змінити його, виконують над похідними представленнями з підтвердженим походженням.

Структурний аналіз виконують над початковою байтовою послідовністю, а піксельні, частотні й часові ознаки формують із контрольовано одержаних похідних даних. До реєстрації рішення об’єкт залишається в ізоляції та не передається до публікації. Це правило зберігає доказове походження структурних відхилень і запобігає використанню неперевіреного носія прикладними компонентами вебсистеми.

Логіку збереження двох джерел свідчень показано на рис.~\ref{fig:process_organization}. Статичний опис $H$ формують із незмінного контейнера та його контрольованих представлень, тоді як початковий журнал подій $J_b$ надходить окремо. Передчасне об’єднання неприпустиме, оскільки $J_b$ до перевірки належності може містити фонові записи, а статичний опис --- спостереження, одержані іншою версією засобу. Метод поведінкового узгодження формує з $J_b$ причинно пов’язану послідовність $X_b$; після перевірки походження статичний опис і ознаки подій і поведінки утворюють інтегроване представлення $Z$, з якого формується аналітичний стан $S_A$. П’ять етапів відповідної технології деталізовано на рис.~\ref{fig:information_technology}.

\begin{figure}[H]
\centering
\begin{tikzpicture}[x=1cm,y=1cm,scale=0.93,transform shape]
\node[dissertation small,text width=2.35cm,minimum height=12mm] (object) at (-6.35,1.25)
  {Вхід 1\\Об’єкт $X_{mm}$};
\node[dissertation small,text width=2.75cm,minimum height=12mm] (staticop) at (-3.25,1.25)
  {Формування $X_c$\\
   Статичний аналіз $F_H$};
\node[dissertation small,text width=1.85cm,minimum height=12mm] (profile) at (-0.55,1.25)
  {Результат методу 1\\Опис $H$};

\node[dissertation small,text width=2.35cm,minimum height=12mm] (log) at (-6.35,-1.25)
  {Вхід 2\\Журнал подій $J_b$};
\node[dissertation small,text width=2.75cm,minimum height=12mm] (events) at (-3.25,-1.25)
  {Причинний відбір\\
   $J_b\!\rightarrow\!X_b$ і кодування $F_b$};
\node[dissertation small,text width=1.85cm,minimum height=12mm] (behavior) at (-0.55,-1.25)
  {Проміжний результат\\Ознаки $Z_b$};

\node[dissertation small,text width=2.25cm,minimum height=18mm] (integration) at (2.05,0)
  {Метод 2\\Контроль походження\\
   та доступності\\Інтеграція $F_I$};
\node[dissertation small,text width=2.21cm,minimum height=12mm] (joint) at (4.85,0)
  {Інтегроване\\представлення $Z$};
\node[dissertation small,text width=2.30cm,minimum height=12mm] (state) at (7.65,0)
  {Модельне рішення $F_D$\\
   Вихід: стан $S_A$};

\node[dissertation container,fit=(object)(staticop)(profile),
  label={[font=\footnotesize]above:Метод 1: формування статичного опису}] {};
\node[dissertation container,fit=(log)(events)(behavior),
  label={[font=\footnotesize]below:Метод 2: причинно-фазовий відбір}] {};

\draw[dissertation arrow] (object.east) -- (staticop.west);
\draw[dissertation arrow] (staticop.east) -- (profile.west);
\draw[dissertation arrow] (log.east) -- (events.west);
\draw[dissertation arrow] (events.east) -- (behavior.west);
\draw[dissertation arrow] (profile.east) -- (integration.155);
\draw[dissertation arrow] (behavior.east) -- (integration.205);
\draw[dissertation arrow] (integration.east) -- (joint.west);
\draw[dissertation arrow] (joint.east) -- (state.west);
\end{tikzpicture}
\caption{Роздільне формування статичного представлення та представлення ознак подій і поведінки з подальшим узгодженням}
\label{fig:process_organization}
\end{figure}


Для зв’язку кожного результату з об’єктом і конкретним етапом його опрацювання процес подано послідовністю станів $S_0,\ldots,S_5$ із визначеним складом. Початковий стан $S_0=\langle X_{mm},C_0\rangle$ містить прийнятий об’єкт і відомості $C_0$ про його надходження. Після попереднього оброблення утворюється $S_1=\langle X_c,m_c^{(1)},C_H\rangle$, де $m_c^{(1)}$ --- початкова маска доступності статичних складових. Результати формування ознак і контрольованого спостереження утворюють $S_2=\langle H,J_b,m_A^{(2)},C_I\rangle$. На цьому етапі $m_A^{(2)}=(m_s,m_v,m_b^{\mathrm{in}})$ містить фактичні значення доступності статичних складових і лише початкове значення доступності журналу; воно не підтверджує причинної придатності подій. Другий метод після перевірки $J_b$ формує повний вхід моделі $X_A$, фактичну маску $m_A^{\mathrm{out}}$ і результат інтеграції $S_3=Z$ або $S_3=Z^{\partial}$. Склад $S_4$ і $S_5$ визначено нижче під час опису прийняття рішення та реєстрації. Кожний наступний стан утворюється з попереднього за параметрами відповідного етапу. Зміна правила розбору, параметра перетворення або порога створює інший результат навіть для незмінного початкового об’єкта. Тому разом зі станом зберігають версії параметрів, а повторне виконання не підміняє раніше одержаного результату.

Для кожного переходу визначено передумову, мінімальний склад результату та стан завершення. Передумова підтверджує належність даних одному об’єкту й сумісність версій. Стан завершення розрізняє повний результат, явно неповний результат і технічну відмову. Відсутність спостереження не свідчить про відсутність аномалії, тому неповний результат передають далі лише разом із маскою доступності та причиною обмеження, а відмову не перетворюють на нульову ознаку.

З урахуванням зазначених умов перехід від прийнятого носія до відтворюваного рішення організовано в такі етапи:

\begin{enumerate}
\item \textit{Попереднє оброблення.} Об’єкту присвоюють ідентифікатор, обчислюють контрольну суму, зберігають незмінний $X_{mm}$, установлюють фактичний формат, формують $X_s$ і $X_m$, а контрольованим розкодуванням одержують $X_v$.
\item \textit{Формування ознак.} Перший метод будує $X_f$ і $X_t$, формує $Z_s$, $Z_v$, оцінки $\rho_s$, $\rho_v$, координати $L_s$, $L_v$ та статичний опис $H$; контрольоване середовище окремо формує початковий журнал $J_b$.
\item \textit{Інтеграція ознак.} Метод поведінкового узгодження відбирає з $J_b$ причинно пов’язані події $X_b$, формує $Z_b$ і виконує перехід
\[
Z=F_I(H,Z_b,m_A,C_I;\theta_I).
\]
Статичні складові надходять до функції інтеграції у складі $H$, а маска та супровідні відомості обмежують допустиме узгодження.
\item \textit{Прийняття рішення.} Функція $F_D$ формує аналітичний стан $S_A$, після чого чинне правило реагування визначає технологічний припис $D_A$ без зміни модельного висновку.
\item \textit{Реєстрація результатів.} Аналітичний стан, технологічний припис, версії, координати, часові межі й посилання на початкові спостереження зберігають в одному записі. Зовнішній припис виконують лише після успішного створення цього запису.
\end{enumerate}

Однакова глибина перевірки всіх об’єктів створює зайві витрати або змушує скорочувати аналіз складних випадків. Підставою для поглиблення є не лише високе значення однієї статичної оцінки, а й суперечність між складовими: структурно нетиповий, але піксельно звичайний об’єкт потребує іншого розбору, ніж об’єкт із двома низькими оцінками. За доступності обох складових їх числову розбіжність характеризує величина $\delta_A=|\rho_s-\rho_v|$, де $\rho_s$ і $\rho_v$ --- відповідно структурна оцінка та оцінка спектральних і піксельних ознак, означені в підрозділі~\ref{sec:svs_method}. До підтвердження спільного калібрування шкал $\delta_A$ використовують лише як діагностичну ознаку конкретної зафіксованої конфігурації, а не як самодостатній доказ змістової суперечності.

Показник $\delta_A$ не визначає загрози, а лише вказує, що дві складові по-різному характеризують стан носія. Передчасне усереднення в такому разі приховало б невизначеність. Спочатку перевіряють повноту обов’язкових даних. Якщо структурна складова недоступна або чинне правило реагування потребує відсутнього піксельного чи кадрового подання, об’єкт передають на експертний розгляд.

За повних даних поглиблений режим перевірки призначають, коли принаймні одна оцінка досягає порога $\tau_r$ або розбіжність досягає порога $\tau_d$; в інших випадках застосовують базовий режим. Базовий режим охоплює визначені статичні ознаки. Поглиблений режим додатково зберігає докладніше частотне представлення позначених ділянок, відомості про розкодування та передбачає контрольоване спостереження оброблення з формуванням $J_b$. Значення $\delta_A$ й обраний режим уносять до супровідних відомостей $C_I$ та враховують під час обчислення упевненості $c_A$.

Пороги $\tau_r$ і $\tau_d$ визначають лише на даних для налаштування та зберігають у версії конфігурації. У контрольних експериментах розділу~\ref{chap:experiments} усі об’єкти опрацьовано за повним статичним режимом. Отже, зазначені пороги характеризують організацію процесу, але не є окремо перевіреною конфігурацією визначення глибини перевірки.

Інформаційний зв’язок етапів узагальнено в таблиці~\ref{tab:ch3_process_stages}. Таблиця фіксує типи входів і результатів кожного етапу та не повторює внутрішніх процедур методів.

\begin{table}[htbp]
\centering
\caption{Входи, основні операції та результати етапів загального процесу виявлення ознак аномального вмісту}
\label{tab:ch3_process_stages}
\small
\begin{tabularx}{\textwidth}{>{\raggedright\arraybackslash}p{0.19\textwidth}YY}
\toprule
\textbf{Етап} & \textbf{Вхід і основна операція} & \textbf{Результат} \\
\midrule
Попереднє оброблення & $S_0$; ідентифікація, структурний розбір і контрольоване розкодування & $S_1=\langle X_c,m_c^{(1)},C_H\rangle$ \\
Формування ознак & $S_1$ і контрольоване середовище; перший метод формує статичні ознаки, а середовище --- початковий журнал & $S_2=\langle H,J_b,m_A^{(2)},C_I\rangle$ \\
Інтеграція ознак & $S_2$; другий метод виконує причинний відбір, формує $X_b$, $Z_b$, $X_A$ та інтегроване представлення & $S_3=Z$ або $S_3=Z^{\partial}$ \\
Прийняття рішення & $S_3$ і $\mathrm T_D$; модель формує $S_A$, а чинне правило реагування --- $D_A$ & $S_4=\langle S_A,D_A\rangle$ \\
Реєстрація результатів & $S_4$, версії та відомості про походження; цілісна реєстрація до зовнішньої дії & $S_5=F_R(S_4,C_R)$ \\
\bottomrule
\end{tabularx}
\end{table}

Загальна організація визначає умови, за яких результат одного етапу може стати входом наступного. Пакет $X_A$ утворюють усередині переходу $S_2\rightarrow S_3$ після причинного відбору; функція $F_I$ перетворює його ознакову частину на $Z$, а $C_I$ зберігає посилання на складники $X_A$, тому вони не втрачаються після переходу. Стан $S_3$ ще не містить технологічного припису: спочатку модельна функція $F_D$ формує $S_A$, а лише потім правило реагування формує $D_A$. Така організація не утворює додаткової моделі або окремого методу, а встановлює незмінні вимоги, без яких ознаки різного походження не можна обґрунтовано зіставляти. У наступних підрозділах ці вимоги конкретизовано в операціях двох методів.

\section{Метод аналізу спектральних, піксельних і структурних ознак мультимедійних даних}
\label{sec:svs_method}

\subsection{Формальна постановка та алгоритм методу}

Модель розділу~\ref{chap:models_formalization} визначила склад статичного опису, але залишила відкритим питання, як одержати його з контейнера, не втративши ознак, що існують лише до розкодування, і не приписавши структурним відхиленням властивостей піксельного вмісту. Просте об’єднання всіх вимірювань в один вектор цієї проблеми не розв’язує. Додаткова байтова ділянка може не змінювати зображення, а слабка зміна коефіцієнтів може не порушувати правил формату. Тому нульовий результат однієї складової не спростовує спостереження іншої.

Метою методу є формування статичного опису ознак можливого прихованого впливу шляхом узгодженого, але роздільного опрацювання двох представлень того самого носія. Структурна складова аналізує початковий контейнер і службові дані та зберігає байтові межі відхилень. Складова спектральних і піксельних ознак аналізує контрольовано розкодоване представлення, його частотні, багатомасштабні, залишкові й часові характеристики та зберігає просторово-часові координати. Об’єднання виконують лише після формування окремих ознак, оцінок і координат. Метод не встановлює остаточного класу й не визначає технологічної дії, оскільки його результат призначено для подальшого причинного зіставлення.

Вхід методу становлять попередньо опрацьовані статичні дані $X_c=\langle X_{mm},\allowbreak X_v,\allowbreak X_s,\allowbreak X_m\rangle$, початкова маска $m_c=(m_s,m_v)$ та супровідні відомості $C_H$. Події $X_b$ до цього входу не включають, щоб не допустити їх неявного впливу на статичний опис. Представлення $X_f$ і $X_t$ утворюють із $X_v$ усередині методу, тому разом із ними зберігають параметри перетворення та відбору. Виходом є $H=\langle Z_v,\allowbreak Z_s,\allowbreak\rho_v,\allowbreak\rho_s,\allowbreak L_v,\allowbreak L_s,\allowbreak m_c,\allowbreak C_H\rangle$. Відповідність між входом і структурованим результатом установлено формулою~\eqref{eq:ch3_svs_mapping}:
\begin{equation}
H=F_H(X_c,m_c,C_H,\Theta_H),
\label{eq:ch3_svs_mapping}
\end{equation}

де $F_H$ --- функція методу аналізу спектральних, піксельних і структурних ознак; $X_c$ --- сукупність статичних представлень; $m_c$ --- початкова маска; $C_H$ --- супровідні відомості статичного аналізу; $\Theta_H$ --- параметри форматних правил, областей, перетворень, відбору кадрів і оцінювання; $H$ --- статичний опис.

Відображення у формулі~\eqref{eq:ch3_svs_mapping} визначає межу відповідальності методу: зміна будь-якого вхідного представлення або параметра створює інший опис і має бути простежуваною. Передумовами виконання є наявність незмінного $X_{mm}$, його контрольної суми та підтверджене походження похідних представлень. Результат має містити не лише числові оцінки, а й координати, маску доступності та супровідні відомості. Відтворювану послідовність формування результату наведено в алгоритмі~\ref{alg:svs_method}. У ньому відокремлено обчислювальні кроки від умов, за яких результат визнають неповним або виконання припиняють.

\begin{algorithm}[H]
\begingroup
\fontsize{12}{13}\selectfont
\setstretch{0.95}
\setlength{\abovecaptionskip}{0pt}
\setlength{\belowcaptionskip}{2pt}
\caption{Формування статичного опису мультимедійного об’єкта}
\label{alg:svs_method}
\textbf{Вхід:} $X_c=\langle X_{mm},X_v,X_s,X_m\rangle$, початкова маска $m_c^{\mathrm{in}}=(m_s^{\mathrm{in}},m_v^{\mathrm{in}})$, супровідні відомості $C_H$, версована конфігурація $\Theta_H$.\par
\textbf{Вихід:} статичний опис $H$ з фактичною маскою $m_c^{\mathrm{out}}$ і статусом повноти в $C_H$ або зареєстрований стан відмови.\par
\textbf{Передумова:} $X_{mm}$ є незмінним; ідентифікатор об’єкта, контрольна сума та походження похідних представлень підлягають перевірці.\smallskip
\begin{enumerate}[label=\arabic*.,leftmargin=6mm,itemsep=1pt,topsep=2pt]
\item Перевірити ідентифікатор, контрольну суму, версії засобів і зв’язок $X_v$, $X_s$, $X_m$ з початковим $X_{mm}$.
\item Якщо незмінність входу або походження хоча б одного використаного представлення не підтверджено, зареєструвати відмову та припинити виконання без аналітичного висновку.
\item Ініціалізувати фактичну маску: $m_c^{\mathrm{out}}\leftarrow m_c^{\mathrm{in}}$.
\item Установити фактичний формат; зіставити заявлений тип, розпізнавальну послідовність байтів, граматичний розбір і результат розкодування; сформувати вектор форматної узгодженості $u_f$.
\item Якщо $m_s^{\mathrm{in}}=1$, походження $X_s$ і $X_m$ підтверджено та структурну карту можна сформувати, обчислити $Z_s=F_s(X_s,X_m;\theta_s)$, оцінку $\rho_s$ і байтові координати $L_s$. Виявлене порушення форматного правила залишити доступною структурною ознакою, а не підставою для $m_s^{\mathrm{out}}=0$. Інакше встановити $m_s^{\mathrm{out}}=0$, $Z_s=\rho_s=L_s=\bot$ і зберегти причину недоступності.
\item Якщо $m_v^{\mathrm{in}}=1$ і контрольоване розкодування виконано успішно, сформувати $X_f$, локальні, частотні, багатомасштабні та залишкові характеристики; для відео відібрати кадри за $\theta_t$, сформувати $X_t$ і виконати часове агрегування. Інакше встановити $m_v^{\mathrm{out}}=0$, $Z_v=\rho_v=L_v=\bot$ і зберегти причину.
\item Якщо $m_v^{\mathrm{out}}=1$, обчислити $Z_v$, оцінку $\rho_v$ і просторово-часові координати $L_v$; якщо розкодування, відбір або обчислення не відтворюються, установити $m_v^{\mathrm{out}}=0$, $Z_v=\rho_v=L_v=\bot$ і зберегти причину.
\item Якщо $m_s^{\mathrm{out}}=m_v^{\mathrm{out}}=0$, зареєструвати технічну незавершеність і припинити виконання без висновку про відсутність аномалії.
\item Сформувати $H=\langle Z_v,Z_s,\rho_v,\rho_s,L_v,L_s,m_c^{\mathrm{out}},C_H\rangle$, не заміщаючи недоступну складову результатом іншої.
\item За потреби обчислити $Z_c=F_C(Z_v,Z_s,m_c^{\mathrm{out}};\theta_C)$ за правилом~\eqref{eq:masked_static_fusion}. Якщо для фактичної маски конфігурацію $\theta_C$ не визначено, установити $Z_c=\bot$ без втрати доступних складників $H$.
\item Повернути $H$ зі статусом «повний», якщо $m_c^{\mathrm{out}}=(1,1)$, або «неповний», якщо $m_c^{\mathrm{out}}\in\{(1,0),(0,1)\}$.
\end{enumerate}
\endgroup
\end{algorithm}

Перші два кроки алгоритму є спільною умовою допустимості подальших обчислень. Маска $m_c^{\mathrm{in}}$ описує очікувану доступність даних до початку обчислень, а $m_c^{\mathrm{out}}$ --- фактично одержані складові; саме її зберігають у $H$. Після контролю входу складову спектральних і піксельних ознак та структурну складову можна обчислювати незалежно; результат однієї складової не заповнює відсутнього результату іншої. Локальний чотирипроцесний запуск, описаний у розділі~4, стосується паралельного опрацювання незалежних об’єктів. Прискорення від паралельного обчислення гілок одного об’єкта окремо не оцінювали.

Потік даних методу показано на рис.~\ref{fig:svs_method}. Суцільні зв’язки утворюють дослідну детерміновану конфігурацію: дві статичні складові формують структурований опис $H$, з якого за потреби одержують числову проєкцію $Z_c$. Пунктирний шлях позначає можливе розширення на основі навченої моделі; він не входить ні до підтвердженого складу $H$, ні до експериментальної проєкції $Z_c$.

\begin{figure}[htbp]
\centering
\begin{tikzpicture}[x=1cm,y=1cm]
\node[dissertation block,text width=3.30cm] (input) at (0,3.15)
  {Статичні дані\\$X_c$};
\node[dissertation block,text width=3.70cm] (structural) at (-4.30,0.80)
  {Структурні ознаки\\$Z_s,L_s$};
\node[dissertation block,text width=3.70cm] (visual) at (0,0.80)
  {Детерміновані спектрально-зорові ознаки\\$Z_v,L_v$};
\node[dissertation block,dashed,text width=3.70cm] (learned) at (4.30,0.80)
  {Неперевірене розширення на основі навченої моделі};
\node[dissertation block,text width=8.10cm] (profile) at (0,-2.10)
  {Статичний опис\\$H=\langle Z_v,Z_s,\rho_v,\rho_s,L_v,L_s,m_c,C_H\rangle$};

\draw[dissertation arrow] (input.south) -- ++(0,-0.40) -| (structural.north);
\draw[dissertation arrow] (input.south) -- (visual.north);
\draw[dissertation arrow,dashed] (input.south) -- ++(0,-0.40) -| (learned.north);
\draw[dissertation arrow] (structural.south) -- ++(0,-0.45) -| (profile.north);
\draw[dissertation arrow] (visual.south) -- (profile.north);
\draw[dissertation arrow,dashed] (learned.south) -- ++(0,-0.45) -| (profile.north);
\end{tikzpicture}
\caption{Структура методу спектрально-зорового та структурного аналізу фото- і відеооб’єктів вебресурсів}
\label{fig:svs_method}
\end{figure}


\subsection{Формування структурних ознак і координат}

Першим етапом є встановлення представлення, яке дозволено аналізувати за правилами формату. Розширення, заявлений тип, розпізнавальна байтова послідовність і результат роботи засобу розкодування можуть збігатися або суперечити одне одному. Сама суперечність ще не доводить прихованого впливу. Окремо перевіряють відповідність розширення заявленому типу, заявленого типу --- розпізнавальній послідовності, цієї послідовності --- установленому під час розкодування типу, а також успішність розкодування. Результати перевірок утворюють вектор $u_f$.

Роздільне збереження умов потрібне для пояснення причини визначення глибини подальшої перевірки. Наприклад, помилка розкодування й наявність даних після логічного завершення контейнера потребують різних перевірок, хоча обидві можуть підвищити загальну структурну оцінку. Сам вектор $u_f$ не визначає прихованого впливу; він указує, де виникла неузгодженість і яку частину структури необхідно дослідити докладніше.

Після встановлення фактичного формату потрібно перейти від переліку контейнерних елементів до ознак, які можна порівнювати між об’єктами, не втрачаючи службових значень. Цей перехід конкретизує модельну функцію $F_s$ у формулі~\eqref{eq:ch3_structural_profile}:
\begin{equation}
Z_s=F_s(X_s,X_m;\theta_s),
\label{eq:ch3_structural_profile}
\end{equation}

де $Z_s$ --- структурне представлення ознак; $X_s$ --- карта контейнера; $X_m$ --- представлення службових даних; $\theta_s$ --- версовані структурні правила та порядок компонентів.

Функція $F_s$ не виправляє карти контейнера й не відкидає порушених службових значень. Вона переводить їх у зіставний опис, з якого надалі формується $\rho_s$, тоді як початкові байтові координати зберігаються в $L_s$. Отже, $Z_s$ впливає на глибину подальшої перевірки, а зв’язок із $X_s$ та $L_s$ дає можливість пояснити цей вплив.

Відтворюваність структурного аналізу залежить від однакового опису звичайних, вкладених і порушених елементів. Для кожного елемента $s_k$ зберігають його тип $c_k$, напіввідкритий байтовий проміжок $[a_k,b_k)$, посилання $w_k$ на батьківський елемент і результат перевірки $v_k$. Права межа $b_k$ до проміжку не входить, а довжина елемента дорівнює $b_k-a_k$. Такий склад є правилом реєстрації: без меж неможливо відтворити координати, без батьківського елемента --- перевірити вкладеність, а без результату перевірки --- відрізнити прочитаний елемент від коректного.

Бінарний результат перевірки $v_k\in\{0,1\}$ набуває одиничного значення лише тоді, коли межі задовольняють умову $0\leq a_k<b_k\leq n$ і виконано форматні обмеження для відповідного типу, де $n$ --- довжина $X_{mm}$ у байтах. Поділ на дві умови необхідний, оскільки фізично доступний інтервал може мати неправильний порядок або недопустиму вкладеність. Якщо заявлена довжина виходить за межі об’єкта, метод зберігає доступну частину, початкове значення поля та причину порушення. Байти після логічного завершення контейнера реєструють як окрему ділянку, а не відкидають і не приписують останньому елементу. Подальший висновок унаслідок цього спирається на фактично прочитані дані, а не на виправлене тлумачення контейнера.

\begin{table}[htbp]
\centering
\caption{Структурні перевірки мультимедійних форматів та очікувані результати розбору}
\label{tab:ch3_format_checks}
\small
\begin{tabularx}{\textwidth}{>{\raggedright\arraybackslash}p{0.12\textwidth}YYY}
\toprule
\textbf{Формат} & \textbf{Контрольовані структури} & \textbf{Ознаки неузгодженості} & \textbf{Результат} \\
\midrule
JPEG & Початок і кінець, службові сегменти, таблиці, послідовності сканування & Порушені межі або порядок, дані після завершення & Карта сегментів і невіднесених ділянок \\
PNG & Розпізнавальна послідовність байтів, основний заголовок, блоки даних, завершення, контрольні суми & Неможлива довжина, помилка контрольної суми, недопустиме повторення & Карта блоків і ознаки цілісності \\
WebP & Контейнер RIFF (формат обміну ресурсними даними), блоки зображення, службові та анімаційні блоки & Розбіжність довжин, несумісний порядок, неврахований залишок & Карта контейнера й альтернативних тлумачень \\
Відео\-контей\-нери & Ієрархічні елементи, доріжки, часові шкали, індекси та потоки & Некоректна вкладеність, суперечлива тривалість, розрив часових міток & Ієрархічна карта й часові ділянки відхилень \\
\bottomrule
\end{tabularx}
\end{table}

Таблиця~\ref{tab:ch3_format_checks} не означає однакової повноти реалізації для всіх форматів. За наявними матеріалами контрольний цикл у розділі~\ref{chap:experiments} обмежено об’єктами у форматах PNG і MP4. Для форматів JPEG і WebP визначено правила розширення, але їх підтримку не слід вважати експериментально підтвердженою без відповідних аналізаторів і окремої перевірки.

Окремі формати містять різну кількість елементів і правил, тому абсолютна кількість спрацювань не є порівнюваною: одна помилка серед двох перевірок і одна серед ста мають різний зміст. Тому порушення спочатку об’єднують у сім змістових груп: межі, порядок, службові дані, невіднесені ділянки, вкладені структури, одночасне тлумачення за правилами кількох форматів та узгодженість розкодування. Для кожної групи $g=1,\ldots,7$ обчислюють частку порушених умов $d_g=n_g^{-}/n_g$, де $n_g^{-}$ --- кількість порушень, а $n_g>0$ --- кількість фактично виконаних перевірок. Нормування приводить частки до спільного числового діапазону, але не доводить змістової порівнянності різних форматів; таку порівнянність потрібно перевіряти окремо.

Самої частки недостатньо, оскільки нуль може означати як відсутність порушень, так і неможливість виконати перевірку. Тому структурне представлення $Z_s$ містить вектор форматних розбіжностей $u_f$, сім часток $d_1,\ldots,d_7$, вектор застосовності $q_s^{(a)}$ і вектор виконання $q_s^{(e)}$. Для групи, застосовної до встановленого формату, відповідна компонента $q_s^{(a)}$ дорівнює одиниці; для незастосовної --- нулю. Компонента $q_s^{(e)}$ дорівнює одиниці лише за фактичного виконання перевірки. Частку $d_g$ обчислюють за $q_{s,g}^{(e)}=1$; за застосовної, але невиконаної перевірки її позначають $\bot$, а не нулем. Унаслідок такого розмежування незастосовність відрізняється від технічної недоступності, невідоме значення не трактується як підтверджена норма, а порушення не змішується з невиконаною перевіркою. Порядок компонентів і версія схеми зберігаються разом з описом.

Повний $Z_s$ потрібний для пояснення висновку, тоді як визначення глибини подальшої перевірки потребує однієї зіставної величини. Після формування вектора обчислюють структурну оцінку $\rho_s=\sigma(\beta_0+\beta^{\mathsf T}Z_s)$, де $\sigma(x)=1/(1+e^{-x})$ --- логістична функція зі значеннями в інтервалі $(0,1)$, а параметри $\beta_0$ і $\beta$ установлено на даних для налаштування. Без окремого калібрування ця оцінка не має імовірнісного змісту. Разом із нею зберігають ознаки та відповідні байтові координати $L_s$, за якими встановлюють порушення, що змінило режим аналізу.

\subsection{Формування спектральних і піксельних ознак}

Складова спектральних і піксельних ознак враховує інше обмеження: слабке місцеве відхилення може зникнути під час усереднення всього зображення або кадру. Тому $X_v$ поділяють на області однакового розміру та кроку, зафіксовані в $\Theta_H$. Для кожної області $r$ окремо визначають частку одиниць у молодшій бітовій площині $b_r$, ентропію Шеннона локальної гістограми значень $e_r$, характеристику переходів між сусідніми значеннями $p_r$, міжканальну різницю $c_r$ і середній модуль високочастотного залишку $h_r$. Ці величини утворюють локальний опис $g_r$. Точні означення статистик $p_r$, $c_r$ і $h_r$, використані фільтри, канали та правила нормування фіксують у складі $\Theta_H$, а їхню версію вносять до супровідних відомостей $C$.

Локальний опис не є самостійним рішенням щодо прихованого навантаження. Текстура, різкий край або сенсорний шум можуть створити подібний відгук без аномального коду. За цим описом визначають ділянки, для яких потрібно зберегти докладніше частотне представлення та відомості про розкодування. Остаточне тлумачення виконують лише після зіставлення з іншими областями, структурними ознаками та доступними подіями.

Локальні статистики не розкривають, у якому масштабі або частотному діапазоні виникло відхилення. Тому піксельне або кадрове подання $X_v$ додатково перетворюють на частотні дані $X_f$. Для цього окремо застосовують дискретне косинусне перетворення, дискретне багатомасштабне перетворення і залишкові фільтри; їхні рівні розкладу, частотні області та параметри утворюють кортеж $\theta_f$, який входить до $\Theta_H$. Роздільне збереження каналів необхідне через різне походження їхніх значень: підсумовування на цьому етапі не дозволило б установити, яке саме перетворення створило відгук.

Під час аналізу об’єктів у форматі JPEG початкові квантовані коефіцієнти, якщо вони доступні, зберігають окремо від коефіцієнтів, повторно обчислених після розкодування. Це обмеження усуває помилкове зіставлення величин, одержаних з різних станів об’єкта. Для кожної області та кожного частотного або залишкового каналу визначають середній модуль коефіцієнтів, стандартне відхилення й ентропію гістограми. Їхній упорядкований набір утворює опис $u_r$.

Для відмежування відтворюваної конфігурації від загальної схеми склад параметрів $\Theta_H$ наведено в табл.~\ref{tab:ch3_theta_h}. Числове значення вважається встановленим лише тоді, коли воно підтверджене експериментальним реєстром; відсутні значення не замінено припущеннями.

\begin{longtable}{@{}p{0.28\textwidth}p{0.27\textwidth}p{0.35\textwidth}@{}}
\caption{Наявні відомості про параметри $\Theta_H$ та стан їх відтворюваності}\label{tab:ch3_theta_h}\\
\toprule
Група параметрів & Зафіксоване значення або правило & Доказовий статус \\
\midrule
\endfirsthead
\multicolumn{3}{l}{Продовження таблиці \thetable}\\
\toprule
Група параметрів & Зафіксоване значення або правило & Доказовий статус \\
\midrule
\endhead
\bottomrule
\endfoot
Подання зображення & $256\times256$ у внутрішньому PNG-наборі; для зовнішніх відеофрагментів --- $320\times240$ зі збереженням пропорцій доповненням полів & Підтверджено реєстрами описаних серій \\
Кадрове подання & $12$ послідовних кадрів із ділянки на рівні $20$~\% тривалості для зовнішніх фрагментів & Підтверджено протоколом зовнішніх серій \\
Локальне вікно і крок & Переносна еталонна реалізація опрацьовує цілий розкодований кадр; історичні локальні параметри не збережено & Повнокадрову реалізацію відтворено; локальну конфігурацію повного методу не перевірено \\
Гістограма & В еталонній реалізації ентропію обчислюють за 256 рівнями восьмибітового сірого подання & Еталонне правило відтворюється; тотожність історичному правилу не встановлено \\
Залишкові фільтри & В еталонній реалізації застосовано лапласіан із ядром $3\times3$ до нормованих каналів & Еталонне правило відтворюється; історичне ядро й нормування не встановлено \\
Дискретне косинусне й дискретне багатомасштабне перетворення & В еталонній реалізації застосовано повнокадрове косинусне перетворення та однорівневе перетворення Гаара & Еталонне правило відтворюється; історичні межі частотної області не встановлено \\
Локальний поріг $\tau_v$ & Добирається лише на налаштувальній частині та фіксується до оцінювання & Числове значення відсутнє в рукописі \\
Структурна оцінка & Порядок дев’яти ознак, середні значення, масштаби, $\beta_0$, $\beta$ і поріг рішення & Відновлено в машинозчитуваному описі параметрів; числову тотожність архівних оцінок перевірено \\
Оцінка спектральних і піксельних ознак & Порядок шістнадцяти ознак, середні значення, масштаби, $\gamma_0$, $\gamma$ і поріг рішення & Відновлено в машинозчитуваному описі параметрів; числову тотожність архівних оцінок перевірено \\
Спільна статична проєкція & Порядок двадцяти чотирьох ознак, середні значення, масштаби, вільний член, коефіцієнти й поріг рішення & Відновлено в тому самому описі; стосується проєкції $Z_c$, а не заміни складників $H$ \\
Агрегування кадрів & $\omega_j=q_j/\sum_{l=1}^{n_v}q_l$ для кадрів із додатною сумою показників технічної якості & Правило не має настроюваних коефіцієнтів; у канонічних таблицях покадрові $q_j$ не збережено \\
Об’єднання часових ділянок & Граничний проміжок між сусідніми кадрами однієї сцени & Числове значення в рукописі не збережено \\
\end{longtable}

Таблиця не є повною конфігурацією запуску. Вона показує, які параметри історичного експерименту можна відтворити з рукопису та реєстрів, а які потребують первинного машинозчитуваного опису. Схема з 12 послідовних кадрів на рівні 20~\% тривалості є окремою скороченою експериментальною конфігурацією. Загальне правило чотирикомпонентного відбору множини $J$, описане нижче, кількісно не перевірено.

Параметри структурної, спектральної та спільної статичної проєкцій збережено в машинозчитуваному описі \nolinkurl{data/model_parameters.json} експериментального пакета. Його контрольна сума SHA-256 дорівнює \nolinkurl{471cc7257d4d7d42b8b97069da82c15d592de84b5d03c39013834a918f4ac4b4}. Числовий вміст відновлено з 29 контрольних записів, оскільки початкового машинозчитуваного файла параметрів у переданому архіві не було. Незалежне застосування відновлених середніх значень, масштабів, коефіцієнтів і порогів до канонічних таблиць ознак повторило 42~168 архівних оцінок трьох проєкцій із максимальною абсолютною похибкою меншою за $10^{-12}$ без зміни жодного порогового рішення. Отже, відтворюваність переходу від збережених ознак до оцінок підтверджено, але побайтову тотожність відновленого опису відсутньому початковому файлу параметрів не стверджено.

До пакета долучено переносну еталонну реалізацію та канонічні таблиці ознак, за якими відтворюються модельні оцінки й підсумкові показники. Однак первинний модуль виділення спектральних ознак не збережено, тому еталонну реалізацію не ототожнено з історичною. Зокрема, не встановлено первинних меж високочастотної області дискретного косинусного перетворення та точного правила обчислення міжканального залишку. Через це числові результати слід тлумачити як результати зафіксованих експериментальних серій, а не як доказ побайтового відтворення повного первинного контуру від початкового медіаоб’єкта до рішення.

Переносна еталонна реалізація з версією від 31 липня 2026 року конкретизує операції, які можна повторити незалежно від архівних ознак. Для кожного розкодованого кадру вона обчислює ентропію за гістограмою з 256 рівнями яскравості, середній модуль і стандартне відхилення лапласіана з ядром $3\times3$, характеристики молодших бітів сірого й синього каналів, міжканальний залишок і часові різниці. Для перетворення Гаара парні розміри кадру поділяють на блоки $2\times2$. Якщо нормовані значення такого блока позначити $a$, $b$, $c$ і $d$, то детальні складові та їхню середню енергію визначено формулою~\eqref{eq:ch3_reference_haar}:
\begin{equation}
\begin{gathered}
H_r=\frac{a+b-c-d}{4},\qquad V_r=\frac{a-b+c-d}{4},\\
D_r=\frac{a-b-c+d}{4},\qquad E_H=\frac{1}{N_r}\sum_{r=1}^{N_r}\left(H_r^2+V_r^2+D_r^2\right),
\end{gathered}
\label{eq:ch3_reference_haar}
\end{equation}

де $H_r$, $V_r$ і $D_r$ --- горизонтальна, вертикальна та діагональна детальні складові блока $r$; $N_r$ --- кількість блоків; $E_H$ --- середня енергія деталей кадру.

Для косинусного перетворення обчислюють матрицю коефіцієнтів $K$ усього сірого або синього каналу. Низькочастотну прямокутну область задають розмірами $h_0=\max(1,\min(8,\lfloor h/8\rfloor))$ і $w_0=\max(1,\min(8,\lfloor w/8\rfloor))$, де $h$ та $w$ --- висота й ширина кадру. Частку енергії поза цією областю визначено формулою~\eqref{eq:ch3_reference_dct}:
\begin{equation}
R_K=\frac{\sum_{u,v}K_{uv}^2-\sum_{u<h_0,\,v<w_0}K_{uv}^2}{\max\!\left(\sum_{u,v}K_{uv}^2,\varepsilon\right)},
\label{eq:ch3_reference_dct}
\end{equation}

де $R_K$ --- еталонна частка високочастотної енергії; $\varepsilon$ --- додатна машинна стала, що запобігає діленню на нуль. Ці правила визначають відтворювану переносну конфігурацію, але не відновлюють невідомих операцій первинного історичного засобу виділення ознак.

Під час узагальнення потрібно одночасно врахувати поширеність слабкого відхилення й одиничний сильний прояв. Використання лише середнього послабило б коротку зміну, а використання лише найбільшого значення зробило б опис надмірно чутливим до випадкового краю або шуму. Тому детерміноване представлення $z_d$ поєднує покомпонентні середні та найбільші значення локальних описів $g_r$ і частотних описів $u_r$. У такий спосіб зберігаються дві властивості відхилення: поширеність і місцева виразність. Водночас $z_d$ залишається описом ознак, а не висновком про наявність аномального коду.

Можливе розширення на основі навченої моделі повинно додатково зберігати положення ділянки, рівень частотного розкладу та формат носія, щоб ураховувати зв’язки між віддаленими областями. Для цього локальні описи можна впорядкувати за просторовими та масштабними індексами й передати попередньо навченій моделі кодування. За наявними матеріалами таку модель не побудовано й не перевірено незалежно. Тому в дослідній конфігурації використано лише $z_d$, а відсутній результат розширення не замінено нульовим вектором. Формальний опис можливого компонента за цих умов не набуває статусу фактично одержаного результату.

Після локального та високочастотного аналізу потрібно утворити вихід складової спектральних і піксельних ознак з визначеною структурою, не приховуючи способу його одержання. Цей перехід конкретизує функцію $F_v$ у формулі~\eqref{eq:ch3_spectral_encoding}:
\begin{equation}
Z_v=F_v(X_v,X_f,X_t;\theta_v),
\label{eq:ch3_spectral_encoding}
\end{equation}

де $Z_v$ --- подання спектральних і піксельних ознак; $X_v$ --- контрольовано одержане піксельне або кадрове подання; $X_f$ --- частотне подання; $X_t$ --- часове подання кадрів; $F_v$ --- функція формування ознак; $\theta_v$ --- версовані параметри формування спектральних і піксельних ознак.

У повній конфігурації методу $F_v$ використовує детерміновані локальні характеристики, характеристики дискретного косинусного й дискретного багатомасштабного перетворень та залишкові характеристики. Додаткове представлення, сформоване навченою моделлю, може ввійти до $Z_v$ лише після окремого навчання, фіксації параметрів і незалежної перевірки. Формула~\eqref{eq:ch3_spectral_encoding} визначає спільний вихід методу, зберігаючи відмінність між фактично доступною і можливою розширеною конфігураціями.

У зафіксованих експериментальних серіях використано скорочену об’єктну конфігурацію: спочатку покадрові характеристики узагальнювали середнім або найбільшим значенням, утворюючи шістнадцятикомпонентний вектор $Z_v$, а потім один раз застосовували логістичне відображення. Канонічні таблиці не містять покадрових оцінок $\rho_{v,j}$, показників якості $q_j$ або координат $L_v$. Отже, вони підтверджують властивості зафіксованої об’єктної проєкції, але не є кількісною перевіркою локального поділу, покадрового зважування чи просторово-часової локалізації повної конфігурації методу.

Як і для структурної складової, повний $Z_v$ зберігають для пояснення, а для визначення глибини перевірки обчислюють проміжну оцінку. Для кожного опрацьованого кадру $j$ її задає логістичне перетворення $\rho_{v,j}=\sigma(\gamma_0+\gamma^{\mathsf T}Z_{v,j})$ з параметрами $\gamma_0$ і $\gamma$, налаштованими окремо від даних оцінювання. Для зображення, яке містить єдиний кадр, $\rho_v=\rho_{v,1}$, а для відео покадрові оцінки узагальнюють за формулою~\eqref{eq:ch3_frame_aggregation}. Покомпонентні значення $Z_v$ зберігають разом із просторовими областями, оскільки сильний високочастотний відгук може бути наслідком текстури, межі, сенсорного шуму або повторного стискання. Тому $\rho_v$ впливає на глибину перевірки, але не використовується як самостійний висновок без урахування формату та відповідної ділянки.

Відмінність між фактично описаним детермінованим формуванням ознак і необов’язковим розширенням на основі навченої моделі показано на рис.~\ref{fig:wavelet_tokenization}. Суцільний шлях завершується підтвердженим виходом $Z_v=z_d$ із координатами $L_v$. Пунктирний шлях завершується на непобудованій моделі кодування і не приєднується до $Z_v$, тому позначає лише напрям подальшого розвитку, а не виконаний експеримент.

\begin{figure}[htbp]
\centering
\begin{tikzpicture}[x=1cm,y=1cm]
\node[dissertation block,text width=3.20cm] (visual) at (-5.80,3.00)
  {Вхід за $m_v=1$\\
   $X_v,X_t,\theta_v$};
\node[dissertation block,text width=4.00cm] (transforms) at (-1.80,3.00)
  {Локальні характеристики, дискретне косинусне й хвилькове перетворення, залишкові фільтри\\
   $X_f$};
\node[dissertation block,text width=3.60cm] (local) at (2.45,3.00)
  {Локальні описи з просторово-часовими координатами};

\node[dissertation block,text width=3.60cm] (deterministic) at (-0.40,0.20)
  {Покомпонентні середні та найбільші значення\\
   $z_d$};
\node[dissertation block,text width=3.80cm] (output) at (-0.40,-2.40)
  {Підтверджена конфігурація\\
   $Z_v=z_d,\;L_v$};
\node[dissertation block,dashed,text width=3.70cm] (sequence) at (4.35,0.20)
  {Упорядкування локальних описів за простором, масштабом і часом};
\node[dissertation block,dashed,text width=3.90cm] (learned) at (4.35,-2.40)
  {Модель кодування не побудовано й не перевірено незалежно};

\draw[dissertation arrow] (visual.east) -- (transforms.west);
\draw[dissertation arrow] (transforms.east) -- (local.west);
\draw[dissertation arrow] (local.240) -- (deterministic.north);
\draw[dissertation arrow,dashed] (local.300) -- (sequence.north);
\draw[dissertation arrow] (deterministic.south) -- (output.north);
\draw[dissertation arrow,dashed] (sequence.south) -- (learned.north);
\end{tikzpicture}
\caption{Формування спектрально-зорових ознак у підтвердженій детермінованій конфігурації та межа неперевіреного розширення}
\label{fig:wavelet_tokenization}
\end{figure}


\subsection{Кадрове агрегування та просторово-часова локалізація}

Для відео аналіз усіх кадрів збільшує витрати й породжує багато залежних спостережень, а аналіз лише ключових кадрів може пропустити короткий прояв. Тому відбір поєднує чотири підстави: належність до ключових кадрів, рівномірне покриття часу, межу сцени та близькість до структурно визначеної ділянки потоку. Об’єднання цих множин утворює $J$ без дублікатів. Крок рівномірного покриття, спосіб визначення меж сцен і радіус часової близькості до структурно визначеної ділянки належать до параметрів $\Theta_H$ і фіксуються у версії конфігурації. Такий спосіб зменшує кількість опрацьованих кадрів за визначеним правилом; повноту збереження коротких проявів окремо не встановлено.

Самого номера кадру недостатньо для відтворення перевірки. Тому в часовому представленні $X_t$ для кожного відібраного кадру $f_j$ зберігають часову мітку $t_j$, усі причини відбору $u_j$ і показник технічної якості розкодування $q_j$. Якщо кадр відповідає кільком умовам, він залишається одним записом, а причини накопичуються. Структурна причина впливає лише на включення кадру до $J$ і відомості про його походження; вона не змінює ні покадрової оцінки $\rho_{v,j}$, ні ваги агрегування.

Покадрові спостереження мають неоднакову технічну придатність: оцінка, одержана з пошкодженого або неповністю розкодованого кадру, не повинна мати той самий вплив, що й оцінка за коректного розкодування. Тому вагу визначають лише за показником технічної якості $q_j\in[0,1]$, обчисленим до одержання $\rho_{v,j}$ за версованим правилом перевірки розкодування. У мінімальній відтворюваній конфігурації $q_j=1$, якщо засіб розкодування повернув непорожній кадр очікуваних розмірів і складу каналів без зареєстрованої помилки, та $q_j=0$ в іншому разі. Проміжне значення $q_j\in(0,1)$ дозволено лише за наявності окремо зафіксованої кількісної функції якості; без такого опису його не призначають. Покадрова оцінка, причина відбору та структурна близькість не входять до ваги; отже, високе значення $\rho_{v,j}$ не збільшує власного внеску під час агрегування. Кадр із $q_j=0$ вважають непридатним і не включають до агрегованого опису. Агрегування виконують лише за $n_v\geq1$ і $\sum_{l=1}^{n_v}q_l>0$. Якщо додатного показника не має жоден відібраний кадр, установлюють $m_v=0$ та $Z_v=\rho_v=L_v=\bot$ і реєструють технічну причину недоступності; рівне усереднення непридатних кадрів не застосовують.

Після відбору й нормування для відео формують один векторний опис та одну проміжну оцінку; обидва результати спираються на однакові ваги. Цю залежність визначено формулою~\eqref{eq:ch3_frame_aggregation}:
\begin{equation}
\omega_j=\frac{q_j}{\sum_{l=1}^{n_v}q_l},\qquad
Z_v=\sum_{j=1}^{n_v}\omega_jZ_{v,j},\qquad
\rho_v=\sum_{j=1}^{n_v}\omega_j\rho_{v,j},
\label{eq:ch3_frame_aggregation}
\end{equation}

де $n_v$ --- кількість різних відібраних кадрів; $q_j$ --- показник технічної якості розкодування кадру; $Z_{v,j}$ --- представлення ознак кадру; $\rho_{v,j}$ --- його проміжна оцінка; $\omega_j$ --- нормована вага кадру; $Z_v$ і $\rho_v$ --- відповідно агреговане представлення та оцінка відео.

Застосування спільних ваг у формулі~\eqref{eq:ch3_frame_aggregation} встановлює однакове правило узагальнення для векторного результату та проміжної числової оцінки. Оскільки покадрові оцінки $\rho_{v,j}$ належать інтервалу $(0,1)$, а ваги невід’ємні з одиничною сумою, агрегована оцінка $\rho_v$ також залишається в цьому інтервалі та не виходить за межі найменшої і найбільшої оцінок придатних кадрів. Водночас велика вага кадру не є доказом аномального коду: вона лише визначає його внесок відповідно до зафіксованого правила технічної якості.

Для фото за $n_v=1$ і $q_1>0$ маємо $\omega_1=1$. Якщо всі придатні кадри мають однаковий показник якості, формула зводиться до звичайного середнього. Покадрових значень $q_j$ у канонічних таблицях історичних серій не збережено, а переносна еталонна реалізація рівномірно усереднює успішно розкодовані кадри. Тому перевагу зважування за якістю над рівним, максимальним або іншим способом узагальнення за наявним корпусом експериментально не встановлено.

Агрегована оцінка характеризує стан відео загалом, але не показує часової ділянки, у якій виникла підстава рішення. Тому функція $G_v$ формує множину координат $L_v=G_v(\{Z_{v,j},\rho_{v,j},q_j,t_j\}_{j=1}^{n_v},\tau_v;\Theta_H)$. Спочатку залишають технічно придатні кадри з $q_j>0$ і $\rho_{v,j}\geq\tau_v$, після чого впорядковують їх за часом. Об’єднують лише сусідні за часом кадри однієї сцени; граничний проміжок об’єднання зафіксовано серед параметрів $\Theta_H$. Для кожної групи зберігають початкову й кінцеву часові мітки, сукупність фактично спостережених областей і найбільшу локальну оцінку. Неспостережений проміжок не оголошують підтвердженою ділянкою, оскільки для нього немає початкового спостереження.

\subsection{Властивості, складність і межі перевірки методу}

Етапи методу та їхні виходи узагальнено в таблиці~\ref{tab:ch3_svs_stages}.

\begin{table}[H]
\centering
\caption{Етапи методу аналізу спектральних, піксельних і структурних ознак, їхні операції та результати}
\label{tab:ch3_svs_stages}
\footnotesize
\renewcommand{\arraystretch}{0.96}
\setlength{\tabcolsep}{4pt}
\begin{tabularx}{\textwidth}{@{}>{\raggedright\arraybackslash}p{0.17\textwidth}>{\raggedright\arraybackslash}p{0.50\textwidth}Y@{}}
\toprule
\textbf{Етап} & \textbf{Операції} & \textbf{Результат} \\
\midrule
Контроль входу & Перевірка ідентифікатора, контрольної суми й походження & Узгоджений $X_c$ \\
Форматна перевірка & Зіставлення заявленого типу, розпізнавальної послідовності байтів, граматики й засобу розкодування & $u_f$ \\
Структурна перевірка & Перевірка типів, меж, вкладеності та службових значень & $L_s$ \\
Структурне кодування & Нормування порушень і врахування застосовності & $Z_s$, $\rho_s$ \\
Локальний аналіз & Молодші біти, ентропія, переходи, міжканальні різниці, залишки & $g_r$ \\
Частотний аналіз & Дискретне косинусне перетворення, дискретне багатомасштабне перетворення, залишкові канали та прив’язка до областей & $X_f$, $z_d$ \\
Оброблення відео & Відбір, формування $X_t$, покадровий аналіз, зважування й визначення часових ділянок & $X_t$, $Z_v$, $\rho_v$, $L_v$ \\
Формування опису & Узгодження двох статичних складових із маскою та супровідними відомостями & $H$ \\
\bottomrule
\end{tabularx}
\end{table}

Обчислювальну складність детермінованої конфігурації методу визначено відносно обсягу фактично опрацьованих представлень. Нехай $n_b=n$ --- кількість байтів початкового мультимедійного об’єкта, $n_p$ --- сумарна кількість піксельних позицій у відібраних кадрах, $n_e=n_s$ --- кількість структурних елементів контейнера, а $n_z$ --- кількість ознак, переданих до обчислення часткових оцінок. За послідовного розбору контейнера структурна складова має порядок зростання часових витрат $O(n_b)$ і потребує до $O(n_e)$ пам’яті для структурної карти. Локальні операції з фіксованими розмірами блоків мають порядок $O(n_p)$. Однак в еталонній конфігурації окремо обчислюють повнокадрове дискретне косинусне перетворення, тому загальна верхня оцінка для складової спектральних і піксельних ознак дорівнює $O(n_p\log n_p)$. Верхня оцінка пам’яті дорівнює $O(n_p)$, якщо розкодовані кадри зберігаються одночасно. Формування спільного статичного опису має порядок $O(n_b+n_p\log n_p)$, а застосування лінійних правил обчислення часткових оцінок --- $O(n_z)$. Ці оцінки характеризують залежність від розміру входу, але не враховують сталих витрат конкретного засобу кодування, розкодування та програмної реалізації. Фактичні часові витрати виміряно окремо в підрозділі~\ref{sec:ch4_ind_svs}.

Результат методу визначається не лише значеннями $\rho_s$ і $\rho_v$. За доступності обох складових утворюється повний статичний опис, а за доступності однієї --- неповний. Відмову фіксують, якщо не підтверджено незмінності входу або походження даних. Відмова й неповнота не перетворюються на висновок про безпечність, оскільки в такому разі відсутня підстава для частини тверджень.

Удосконалення методу полягає в узгодженому роздільному опрацюванні двох статичних складових, збереженні байтових і просторово-часових координат та формуванні одного статичного опису без втрати відомостей про походження. Одержаний $H$ є необхідним проміжним результатом: він уможливлює зіставлення статичного спостереження з подією, але сам не доводить прояву прихованого впливу. Межі методу визначаються переліком форматів, для яких реалізовано правила аналізу, якістю розкодування, параметрами кадрового відбору й відсутністю окремо перевіреного засобу кодування піксельних ознак на основі навченої моделі. Групово відокремлену перевірку скорочених числових проєкцій $Z_s$, $Z_v$ і $Z_c$ на об’єктах у форматах PNG і MP4 наведено в підрозділі~\ref{sec:ch4_ind_svs}; локалізації, покадрове зважування та повний склад $H$ у цій серії окремо не оцінювали. Потреба встановити зв’язок між статичним описом і реакцією середовища зумовлює перехід до другого методу.

\section{Метод поведінкового узгодження ознак різного походження}
\label{sec:behavior_alignment}

Статичний опис установлює, де й у якому представленні спостерігається відхилення, але не показує його зв’язку з реакцією компонента вебсистеми. Журнал оброблення, навпаки, містить дії процесів, однак без перевірки походження не дозволяє відрізнити наслідок опрацювання поточного об’єкта від фонової діяльності. Безпосереднє приєднання журналу до $H$ створило б хибний причинний зв’язок, а відмова від урахування подій залишила б невстановленим можливий прояв прихованого впливу.

Межі підтверджених у підрозділі результатів визначено наявними матеріалами. Доказовий статус розмежовано на три рівні: формальну специфікацію повного методу, експериментально перевірену скорочену модель узгодження статичних оцінок і неперевірену повну конфігурацію з причинно пов’язаним журналом. Нижче формалізовано причинно-фазове правило допуску подій, їх числове подання та інтеграцію зі статичним описом; кількісну перевагу повної конфігурації не заявлено. Скорочена модель, адаптована до формату й навчена за експериментальними даними, яку досліджено в розділі~\ref{chap:experiments}, узгоджує вже сформовані статичні оцінки першого методу й не є експериментальною перевіркою повного методу поведінкового узгодження. Зіставлення з найближчими аналогами наведено в табл.~\ref{tab:ch3_novelty_comparison}.

Метою методу є інтеграція статичних ознак із подіями оброблення того самого мультимедійного об’єкта за підтвердженої належності та фазової сумісності. Перевірка цих умов відокремлює просту часову одночасність від змістової підтримки статичного припущення. Метод одержує початковий журнал $J_b$, установлює причинну належність його записів, формує $X_b$ і $Z_b$ та лише після цього виконує інтеграцію. За відсутності придатного журналу статичний опис зберігається як неповний результат, а не доповнюється штучно утвореною нульовою поведінкою.

\subsection{Формальна постановка та етапи методу}

Об’єктом опрацювання методу є пара взаємопов’язаних, але початково розділених джерел спостережень: статичний опис мультимедійного об’єкта $H$ і журнал подій контрольованого оброблення $J_b$. На вході також мають бути доступними ідентифікатор об’єкта, ідентифікатор запуску, часові межі фаз, початкова маска доступності, версії схем ознак, словників і каталогу правил. Ці відомості становлять формальні вимоги до входу. Відсутність обов’язкової відомості не компенсується числовим значенням, оскільки невідоме походження журналу унеможливлює причинне зіставлення.

Формальний опис виходу передбачає один із трьох станів. Повний результат $Z$ містить статичні ознаки й ознаки подій і поведінки, їх координати, фактичну маску доступності, показники повноти, версії конфігурацій та посилання на початкові спостереження. Неповний результат $Z^{\partial}$ зберігає доступні статичні складові й документовану причину відсутності поведінкового внеску. Стан відмови формується тоді, коли не підтверджено цілісність $H$, походження $J_b$ або сумісність версій. Отже, вихід методу не зводиться до одного числового показника, а має формально визначений склад, повноту та простежуваність.

Послідовність перетворень від вхідних даних до одного з вихідних станів подано в табл.~\ref{tab:ch3_behavior_stages}. Кожний наступний етап використовує лише результат, прийнятий попереднім етапом. Це обмеження виключає формування поведінкової оцінки з неперевірених записів і не дозволяє змішувати події різних запусків.

\begin{table}[htbp]
\centering
\small
\caption{Етапи методу поведінкового узгодження ознак різного походження}
\label{tab:ch3_behavior_stages}
\begin{tabularx}{\textwidth}{@{}>{\raggedright\arraybackslash}p{0.19\textwidth}Y>{\raggedright\arraybackslash}p{0.24\textwidth}@{}}
\toprule
Етап & Зміст перетворення & Результат етапу \\
\midrule
Контроль входу & Перевірка цілісності $H$, походження $J_b$, ідентифікаторів, часових меж і сумісності версій & Допущений вхід або стан відмови \\
Причинно-фазовий відбір & Вилучення записів іншого об’єкта, запуску, фази або причинного ланцюга & Послідовність $X_b$ і посилання на початкові записи \\
Унормування подій & Перевірка схеми, упорядкування, кодування дії, ініціатора, ресурсу, результату й часу & Числові описи подій $u_t$ \\
Формування ознак & Застосування версованого каталогу відповідностей до причинно пов’язаної послідовності & Представлення $Z_b$ і оцінка $\rho_b$ \\
Змістове зіставлення & Пов’язування поведінкових компонентів зі статичними координатами та фазами оброблення & Оцінка прояву $Q_A$ і показники повноти \\
Інтеграція та контроль виходу & Формування $Z$ або $Z^{\partial}$, перевірка маски, координат, версій і посилань & Результат методу з формально визначеним складом \\
\bottomrule
\end{tabularx}
\end{table}

Метод не встановлює причинного зв’язку лише за близькістю часових позначок. Час є необхідною, але недостатньою умовою: запис має стосуватися того самого об’єкта або його простежуваної копії, належати тому самому запуску, відповідати дозволеній фазі та мати допустимого причинного попередника. Версія схеми події також має бути сумісною з версією словників, інакше числове кодування одного поля може набути іншого змісту.

Нехай $d_o(e)$ позначає відповідність події об’єкту, $d_r(e)$ --- запуску, $d_{\phi}(e)$ --- фазі, $d_p(e)$ --- причинному попереднику, а $d_v(e)$ --- версії схеми. Складену ознаку допустимості події визначено формулою~\eqref{eq:ch3_event_admissibility}:
\begin{equation}
d(e)=d_o(e)d_r(e)d_{\phi}(e)d_p(e)d_v(e),\qquad d_{\bullet}(e)\in\{0,1\},
\label{eq:ch3_event_admissibility}
\end{equation}

де $d(e)=1$ лише тоді, коли виконано всі п’ять умов. Виконання однієї умови не компенсує порушення іншої. Подія з правильною часовою позначкою, але з іншим номером запуску, не впливає на $Z_b$. Подію поточного об’єкта з невідомою версією схеми також не кодують автоматично.

Правило $d_p(e)$ окремо враховує кореневу подію. Для першої події дозволеного ланцюга $d_p(e)=1$ лише за спеціального значення $\operatorname{parent}(e)=\bot$ і належності фази до заданої в $\theta_b$ множини початкових фаз. Для кожної наступної події посилання має вести до раніше допущеного запису того самого об’єкта й запуску; часова позначка попередника не може бути пізнішою, а послідовність посилань не має утворювати циклу.

Перевірка $d(e)$ потребує мінімального складу початкового запису. Для кожної події зберігають ідентифікатори об’єкта та запуску, фазу, дію, ініціатора, ресурс, результат, часову позначку, посилання на причинного попередника й версію схеми. Якщо обов’язкове поле відсутнє, відповідна складова $d_{\bullet}(e)$ не набуває припущеного значення; запис вилучають з $X_b$, а причину вилучення реєструють. Реєстрація цієї причини уможливлює повторну перевірку відбору без зміни початкового журналу.

Перехід від початкового журналу до причинно пов’язаної послідовності визначено формулою~\eqref{eq:ch3_event_selection}:
\begin{equation}
X_b=P_b(J_b,C_I;\theta_b),
\label{eq:ch3_event_selection}
\end{equation}

де $P_b$ --- функція причинно-фазового відбору; $J_b$ --- незмінний початковий журнал поточного запуску; $C_I$ --- відомості про об’єкт, запуск, фази й походження записів; $\theta_b$ --- версовані правила, словники та пороги поведінкової складової; $X_b$ --- послідовність відібраних причинно пов’язаних подій.

Інтеграції підлягають три групи ознак, а не готові рішення окремих засобів виявлення. Тому вихід методу конкретизує функцію об’єднання моделі у формулі~\eqref{eq:ch3_alignment_mapping}:
\begin{equation}
Z=F_I(H,Z_b,m_A,C_I;\theta_I),
\label{eq:ch3_alignment_mapping}
\end{equation}

де $F_I$ --- функція інтеграції ознак; $H$ --- статичний опис з окремими групами $Z_v$ і $Z_s$, оцінками $\rho_v$, $\rho_s$ та їх координатами; $Z_b$ --- ознаки подій і поведінки; $m_A$ --- маска доступності; $C_I$ --- відомості про походження, якість і параметри інтеграції; $\theta_I$ --- версована конфігурація інтеграції; $Z$ --- інтегроване представлення.

Узгоджені за схемою $Z_v$ і $Z_s$ надходять до формули~\eqref{eq:ch3_alignment_mapping} як окремі складники опису $H$, а не відновлюються з похідного $Z_c$. Маску $m_A$ та відомості $C_I$ передають як окремі аргументи, а координати $L_s$ і $L_v$ містяться в $H$. Перед інтеграцією перевіряють рівності $m_A|_{\{s,v\}}=m_c$ і $C_I|_H=C_H$; за розбіжності інтеграцію припиняють. Тому вихід функції відповідає структурі $Z=\langle z,L,m_A,C_I\rangle$ і зберігає підстави формування інтегрованого числового вектора.

Потік даних, умову допуску журналу та два граничні результати методу показано на рис.~\ref{fig:behavior_alignment}. За придатного журналу метод виконує причинно-фазовий відбір, розрізняє порожню й непорожню спостережувані послідовності, формує $Z_b$ і завершує повну інтеграцію. За недоступного або непридатного журналу поведінкову складову позначають недоступною та формують $Z^{\partial}$ із документованою причиною. Пунктирний блок відображає неперевірене розширення, яке не приєднується до виходу виконуваної конфігурації.

\begin{figure}[htbp]
\centering
\begin{tikzpicture}[x=1cm,y=1cm]
\node[dissertation block,text width=3.25cm] (profile) at (-4.80,3.45)
  {Статичний опис\\$H$};
\node[dissertation block,text width=3.25cm] (events) at (0,3.45)
  {Початковий журнал\\$J_b$};
\node[dissertation block,text width=3.25cm] (context) at (4.80,3.45)
  {Доступність і супровідні\\відомості $m_A,C_I$};
\node[dissertation block,text width=3.25cm] (rules) at (-2.60,1.15)
  {Причинний відбір і явні\\правила $X_b,Z_r$};
\node[dissertation block,dashed,text width=3.25cm] (states) at (2.60,1.15)
  {Неперевірені послідовнісні стани\\$U_b$};
\node[dissertation block,dashed,text width=3.65cm] (attention) at (3.00,-1.05)
  {Неперевірене перехресне зіставлення $A$; до дослідного виходу не входить};
\node[dissertation block,text width=3.25cm] (integration) at (0,-2.85)
  {Інтеграція ознак\\$F_I$};
\node[dissertation block,text width=3.25cm] (joint) at (0,-4.65)
  {Інтегроване представлення\\$Z$};

\draw[dissertation arrow] ([xshift=-0.65cm]events.south)
  -- ++(0,-0.45) -| (rules.north);
\draw[dissertation arrow,dashed] ([xshift=0.65cm]events.south)
  -- ++(0,-0.45) -| (states.north);
\draw[dissertation arrow,dashed] (states.south) -- (attention.north);
\draw[dissertation arrow] (profile.south) -- ++(0,-0.45) -| (integration.west);
\draw[dissertation arrow,dashed] (context.east) -- ++(0.45,0) |- (attention.east);
\draw[dissertation arrow] (context.west) -- ++(-0.45,0) |- (integration.east);
\draw[dissertation arrow] (rules.south) -- ++(0,-0.45)
  -| ([xshift=-0.65cm]integration.north);
\draw[dissertation arrow] (integration.south) -- (joint.north);
\end{tikzpicture}
\caption{Підтверджена та неперевірена гілки формалізованої процедури поведінкового узгодження}
\label{fig:behavior_alignment}
\end{figure}


\subsection{Алгоритм причинного допуску та кодування подій}

Відтворювану послідовність допуску й узгодження подій наведено в алгоритмі~\ref{alg:behavior_alignment}. Алгоритм описує конфігурацію з визначеними правилами; неперевірений шлях кодування послідовностей у ньому вилучено з виконуваної конфігурації. Таке розділення запобігає тлумаченню формалізації можливого розширення як доказу його фактичної реалізації.

\clearpage

\begin{algorithm}[H]
\begingroup
\setlength{\abovecaptionskip}{0pt}
\setlength{\belowcaptionskip}{2pt}
\caption{Метод поведінкового узгодження статичного опису з подіями оброблення}
\label{alg:behavior_alignment}
\fontsize{12}{13}\selectfont
\setstretch{0.95}
\textbf{Вхід:} статичний опис $H$, початковий журнал $J_b$, початкова маска $m_A^{\mathrm{in}}$, супровідні відомості $C_I$, версовані конфігурації $\theta_b$, $\theta_r$ і $\theta_I$.\par
\textbf{Вихід:} повне представлення $Z$ або неповне представлення $Z^{\partial}$ з фактичною маскою $m_A^{\mathrm{out}}$ і позначенням повноти в $C_I$; за порушення передумов --- зареєстрований стан відмови.\par
\textbf{Передумова:} $H$ одержано за алгоритмом~\ref{alg:svs_method}; за наявності $J_b$ його цілісність і походження підлягають перевірці.\smallskip
\begin{enumerate}[label=\arabic*.,leftmargin=6mm,itemsep=0pt,topsep=1pt]
\item Перевірити цілісність $H$ і сумісність $m_A^{\mathrm{in}}|_{\{s,v\}}=m_c$, $C_I|_H=C_H$ та версій схем ознак.
\item Якщо цілісність або сумісність не підтверджено, зареєструвати відмову й припинити інтеграцію.
\item Ініціалізувати фактичну маску $m_A^{\mathrm{out}}\leftarrow(m_s,m_v,m_b^{\mathrm{in}})$, де $(m_s,m_v)=m_c$ зі статичного опису $H$.
\item Якщо $J_b$ недоступний, має недостатнє фазове покриття або непідтверджене походження, установити $Z_b=\bot$ і $m_A^{\mathrm{out}}=(m_s,m_v,0)$, зберегти причину в $C_I^{\partial}$, сформувати $Z^{\partial}=F_I(H,\bot,m_A^{\mathrm{out}},C_I^{\partial};\theta_I)$ за формулою~\eqref{eq:partial_integrated_representation}, повернути $Z^{\partial}$ зі статусом «неповний» і завершити виконання методу.
\item Відібрати з $J_b$ записи поточного запуску, що стосуються об’єкта або його простежуваної копії. Кореневу подію допустити лише за $\operatorname{parent}(e)=\bot$ і дозволеної початкової фази; кожну наступну --- за посилання на раніше допущеного попередника без часової інверсії та циклу; сформувати $X_b=P_b(J_b,C_I;\theta_b)$.
\item Якщо журнал повністю охоплює обов’язкові фази, але придатних подій немає, установити $X_b=\varnothing$, $T=0$ і $m_A^{\mathrm{out}}=(m_s,m_v,1)$; не ототожнювати цей стан із недоступним журналом.
\item Для кожної $e_t\in X_b$ перевірити схему й часовий порядок, унормувати дію, ініціатора, ресурс, результат і час за версованими словниками; зберегти посилання на початковий запис.
\item Обчислити представлення на основі правил $Z_r=F_r(X_b;\theta_r)$ і для формально визначеної конфігурації з явними правилами прийняти $Z_b=Z_r$; обчислити $\rho_b$ і $Q_A$, після чого встановити $m_{bq}=\mathbf 1\!\left[m_b=1\land F_{bq}(\rho_b,Q_A,C_I)\neq\bot\right]$. Якщо $m_{bq}=0$, зберегти причину недоступності спільної поведінкової оцінки й не передавати невизначений внесок до функції прийняття рішення.
\item Не активувати формування послідовнісних станів $U_b$ та перехресне зіставлення $A$, доки відповідну модель не навчено, не версовано й не перевірено незалежно.
\item Установити $m_A^{\mathrm{out}}=(m_s,m_v,1)$, зіставити $L_s$ і $L_v$ з фазами та ресурсами відібраних подій і сформувати $Z=F_I(H,Z_b,m_A^{\mathrm{out}},C_I;\theta_I)$ без заміщення недоступних даних нульовими значеннями.
\item Перевірити наявність у $Z$ фактичної маски, координат, статусу повноти, версій конфігурації та посилань на $H$ і $J_b$; повернути $Z$.
\end{enumerate}
\endgroup
\end{algorithm}

Перші чотири кроки алгоритму не виділяють поведінкових ознак, а встановлюють допустимість впливу журналу на висновок і фактичну маску результату. Якщо журнал непридатний, метод повертає однозначно визначене $Z^{\partial}$ без поведінкового внеску. Наступні кроки формують подієве представлення, перевіряють його змістову відповідність статичним координатам і завершують повну інтеграцію. Частково унормовані записи з непідтвердженим походженням не передають далі як спостереження.

Подія набуває змісту не лише через тип, а й через місце в причинно-часовому ланцюгу. Після відсікання фонових записів придатні події впорядковують у послідовність $X_b=\langle e_1,\ldots,e_T\rangle$, де $T$ --- кількість подій. Якщо журнал охоплює всі обов’язкові фази, але причинно пов’язаних подій не виявлено, встановлюють $X_b=\varnothing$, $T=0$ і $m_b=1$. Значення $m_b=0$ установлюють лише за недоступності журналу, недостатнього фазового покриття або непідтвердженого походження записів. Перелік обов’язкових фаз відповідає компонентам послідовності оброблення --- прийманню, структурному розбору, розкодуванню, перетворенню та збереженню --- і його фіксують у конфігурації методу разом зі словниками. Таке розмежування не дає технічній неповноті журналу зменшити оцінку ризику й водночас зберігає спостережувану відсутність подій як окремий результат.

Граничні стани журналу та відповідні результати методу систематизовано в табл.~\ref{tab:ch3_behavior_boundary_states}. Вони є взаємовиключними: повністю спостережений порожній журнал не може одночасно мати ознаку недоступності, а журнал із несумісною схемою не можна прийняти як повний лише тому, що він містить записи.

\begin{table}[htbp]
\centering
\small
\caption{Граничні стани методу поведінкового узгодження}
\label{tab:ch3_behavior_boundary_states}
\begin{tabularx}{\textwidth}{@{}>{\raggedright\arraybackslash}p{0.24\textwidth}>{\raggedright\arraybackslash}p{0.22\textwidth}Y@{}}
\toprule
Стан вхідних відомостей & Формальне подання & Результат методу \\
\midrule
Журнал недоступний або не охоплює обов’язкові фази & $m_b=0$, $Z_b=\bot$ & Неповне $Z^{\partial}$ із зареєстрованою причиною; поведінковий внесок не обчислюється \\
Журнал повний, придатних подій немає & $m_b=1$, $X_b=\varnothing$ & Спостережуваний порожній результат; правило для нульових лічильників застосовується без висновку про безпечність \\
Журнал повний, допустимі події наявні & $m_b=1$, $T>0$ & Формуються $Z_b$, $\rho_b$, $Q_A$ та повне $Z$ за сумісності всіх складників \\
Походження або версія схеми не підтверджені & $d_o(e)=0$ або $d_v(e)=0$ для істотних записів & Записи не інтегруються; залежно від покриття формується неповний результат або стан відмови \\
Навчуване розширення не перевірене & $U_b$ і $A$ не допущені до виконуваної конфігурації & Результат визначається лише конфігурацією з явними правилами \\
\bottomrule
\end{tabularx}
\end{table}

Розрізнення цих станів потрібне не лише для правильного обчислення маски. Воно впливає на можливість повторного використання результату: $Z^{\partial}$ може бути доповнене після одержання належного журналу лише як новий версований запуск, тоді як повністю спостережене $X_b=\varnothing$ уже є завершеним поведінковим спостереженням для зафіксованих фаз. Попередній запис при цьому не перезаписують, оскільки поява нових подій або зміна каталогу створює іншу конфігурацію доказів.

Однакова дія в різних фазах або над різними ресурсами може мати різний зміст. Для кожної події $e_t$ зберігають дію $a_t$, ініціатора $i_t$, ресурс $r_t$, результат $o_t$ і час $\tau_t$. Фаза оброблення, ідентифікатор процесу та посилання на причинного попередника входять до супровідних відомостей про подію. Події різних фаз, ініціаторів або результатів не об’єднують. Суміжні однотипні записи можна згортати лише зі збереженням їх кількості, часових меж і посилань на початкові записи. Такий порядок скорочення журналу не руйнує причинного ланцюга.

Запис із визначеним складом придатний для перевірки людиною, але обчислювальне зіставлення потребує числового представлення $u_t$. Його утворює функція $E_b$ за словниками категорій з визначеними версіями, часовими масштабами й масками полів. Тип дії та клас ресурсу кодують цілими індексами відповідних словників, результат --- двійковою ознакою успішності, а час --- нормованим зсувом відносно початку запуску. Відсутнє поле позначають окремою маскою, а не значенням, яке може збігатися з реальною категорією. Словники та порядок полів залишаються незмінними в межах однієї версії методу, оскільки їх зміна змінює зміст кожної компоненти.

Окреме числове подання подій не відображає залежностей між віддаленими фазами. Повна конфігурація може додатково формувати впорядковані послідовнісні стани $U_b$ з описів $u_1,\ldots,u_T$. Проте матеріали дослідження не містять побудованої та незалежно перевіреної моделі кодування причинно пов’язаних послідовностей подій. Отже, послідовнісні стани визначають лише можливе розширення і не впливають на зафіксований результат дослідної конфігурації.

Незалежно від обраної конфігурації вихід поведінкової складової має відповідати єдиній структурі моделі. Перехід від причинно пов’язаної послідовності до такого представлення задано функцією $F_b$ у формулі~\eqref{eq:ch3_behavior_encoding}:
\begin{equation}
Z_b=F_b(X_b;\theta_b),
\label{eq:ch3_behavior_encoding}
\end{equation}

де $F_b$ --- функція формування ознак подій і поведінки; $Z_b$ --- їх представлення; $\theta_b$ --- версовані словники, правила й порядок компонентів.

Спільна структура $Z_b$ забезпечує однаковий формат входу функції інтеграції, але не усуває залежності семантики й калібрування від способу одержання ознак подій і поведінки. Разом із результатом зберігають позначення використаної конфігурації, тому вихід, сформований за визначеними правилами, не можна ототожнювати з результатом неперевіреної моделі кодування. Значення $Z_b$ впливає на інтегроване представлення лише за підтвердженої доступності поведінкової складової.

У повній конфігурації $F_b$ може використовувати $U_b$. У формально визначеній базовій конфігурації скорочений варіант ґрунтується на явних правилах. Він застосовує каталог допустимих і нетипових відповідностей між фазою, дією, ресурсом, результатом і причинним попередником із визначеною версією. Нормовані лічильники та ознаки виконання правил утворюють представлення $Z_r=F_r(X_b;\theta_r)$, де $\theta_r$ позначає версію каталогу.

Для формально описаної конфігурації з явними правилами прийнято $Z_b=Z_r$. Таке рішення зумовлене відсутністю перевіреної моделі кодування й не означає рівноцінності сформованого опису можливому розширенню. Для відтворення потрібні унормований $X_b$, словники, каталог правил, порядок компонентів і пороги. Заміна каталогу створює іншу конфігурацію та має відображатися у версії результату. Оскільки повного каталогу в наявному корпусі не збережено, цю конфігурацію не подають як експериментально виконану.

Можливе розширення на основі навченої моделі має розв’язати інше завдання, ніж конфігурація з явними правилами: зіставити не два вже узагальнені вектори, а конкретну статичну ділянку з причинно сумісною подією. Для цього статичні елементи, сформовані з $Z_v$, $Z_s$, $L_v$ і $L_s$, та послідовнісні стани $U_b$ мають бути зведені до спільної розмірності. Однаковий розмір векторів не є достатньою підставою для їх поєднання, оскільки числова подібність може виникнути між подіями різних запусків або несумісних фаз.

Механізм розширення передбачає причинно-фазову маску $M$. Вона дозволяє зіставляти лише пари, що належать одному об’єкту, одному запуску, дозволеному ланцюгу процесів і сумісним фазам. Для дозволених пар обчислюють нормовані ваги відповідності, а їх результат утворює узгоджене представлення $A$. Якщо статичний елемент не має жодної дозволеної події, відповідний рядок не перетворюють на звичайний нульовий вектор, а позначають недоступним. Ця умова не дозволяє тлумачити відсутність спостереження як підтверджену відсутність поведінкового прояву.

Після зіставлення дозволені представлення можна узагальнити й поєднати зі статичними ознаками з урахуванням маски $m_A$. Внесок поведінкових даних має зменшуватися за неповної послідовності подій і бути відсутнім, коли причинну належність не підтверджено. Можливість поєднання визначає маска, а не умовне числове значення відсутньої складової. За наявними матеріалами послідовнісну модель кодування, причинно-фазове зважування та поєднання на основі навченої моделі не налаштовано й не перевірено. Цей механізм описує умови можливого розширення, але не входить до результатів перевірки дослідної конфігурації.

\subsection{Інтеграція та оцінювання результату методу}

У конфігурації з явними правилами використовують поєднання статичного представлення з урахуванням маски $Z_c=F_C(Z_v,Z_s,m_c;\theta_C)$ з вектором $Z_r$, для якого зберігають походження всіх складників. Початкові $Z_v$ і $Z_s$ при цьому не вилучають. Якщо $m_b=0$, поведінковий внесок не обчислюють, а результат зберігає статичні ознаки та явне позначення неповноти. Якщо $m_b=1$, правила визначають, які поведінкові компоненти можуть підтримати або послабити конкретне статичне спостереження. Результатом є пояснення внеску подій, відтворюване лише за повної фіксації правил і параметрів, але не автоматичне доведення програмного змісту прихованого фрагмента.

Інтегроване представлення зберігає можливість подальшої класифікації, але для пояснення поведінкового внеску потрібна окрема проміжна оцінка $\rho_b$. Її визначено формулою~\eqref{eq:ch3_behavior_score}:
\begin{equation}
\rho_b=F_{\rho b}(Z_b;\theta_{\rho b})\in[0,1],
\label{eq:ch3_behavior_score}
\end{equation}

де $F_{\rho b}$ --- версована функція оцінювання ознак подій і поведінки, а $\theta_{\rho b}$ --- її параметри. Оцінку обчислюють лише за $m_b=1$. Для повністю спостереженого $X_b=\varnothing$ функція використовує окремо зафіксоване правило для нульових лічильників; одержане значення не тлумачать як доказ безпечності. За $m_b=0$ установлюють $\rho_b=\bot$, а підсумковий стан містить ознаку неповноти, а не занижену оцінку.

Висока виразність подій ще не означає прояву через конкретний компонент, оскільки та сама послідовність може бути штатною на іншій фазі або в іншому режимі оброблення. Тому оцінка прояву $Q_A=F_Q(Z_b,R_{web},C_I)$ зіставляє поведінкові ознаки зі складом компонентів $R_{web}$ і причинними відомостями $C_I$. Нехай $\mathrm R_{\mathrm{app}}$ --- множина правил каталогу, фактично застосовних до поточної послідовності оброблення; $I_r\in\{0,1\}$ --- ознака виконання правила нетипової відповідності; $\alpha_r>0$ --- його вага. Для непорожньої $\mathrm R_{\mathrm{app}}$ оцінку визначено формулою~\eqref{eq:ch3_manifestation_score}:
\begin{equation}
Q_A=
\frac{\sum_{r\in\mathrm R_{\mathrm{app}}}\alpha_r I_r}
{\sum_{r\in\mathrm R_{\mathrm{app}}}\alpha_r},
\label{eq:ch3_manifestation_score}
\end{equation}

де $Q_A$ --- зважена частка виконаних правил нетипової відповідності серед застосовних правил. Ваги правил зафіксовано у версії каталогу, тому $Q_A\in[0,1]$. Якщо $\mathrm R_{\mathrm{app}}=\varnothing$, оцінку не обчислюють, позначають $\bot$ і зберігають причину; незастосовні правила не утворюють документованого нульового спостереження.

Одержане $Q_A$ входить до $Z$ і зберігається у супровідних відомостях результату. Під час формування інтегральної оцінки ризику його враховують спільно з $\rho_b$ через функцію $F_{bq}(\rho_b,Q_A,C_I)$, визначену формулою~\eqref{eq:behavior_joint_score} і використану у формулі~\eqref{eq:integral_risk}; окреме додавання цих залежних оцінок не допускається. За наявними матеріалами самостійної експериментальної перевірки $Q_A$ не проведено, тому ця величина визначає склад результату методу, але не підтверджує його кількісної переваги.

Після об’єднання потрібно відокремити числову близькість часткових оцінок від якості підстав для їх зіставлення. Нехай $\mathrm I_A=\{k\in\{v,s,b\}:m_k=1\}$, $n_m=|\mathrm I_A|$, а $\rho_k$ --- відповідна доступна оцінка. За $n_m\geq2$ показник числової близькості визначено формулою~\eqref{eq:ch3_numeric_proximity}:
\begin{equation}
q_I=
1-\frac{2}{n_m(n_m-1)}
\sum_{\substack{r<s\\r,s\in\mathrm I_A}}
|\rho_r-\rho_s|,
\label{eq:ch3_numeric_proximity}
\end{equation}

де $q_I$ --- середня числова близькість усіх пар доступних часткових оцінок. За $n_m<2$ установлюють $q_I=\bot$. До перевірки спільного калібрування шкал $q_I$ тлумачать лише як показник числової близькості, а не як підтвердження змістової узгодженості. Однаково високі або низькі помилкові оцінки також можуть бути близькими. Тому $q_I$ не є підсумковим ризиком, не доводить програмного змісту прихованого фрагмента й не збільшує підсумкову упевненість без окремої емпіричної перевірки.

Окремо функція $F_N$ визначає упевненість інтеграції $c_I=F_N(m_A,q_p,q_t)$, де $q_p\in[0,1]$ --- частка охоплених обов’язкових фаз. За $T>0$ величина $q_t$ дорівнює частці подій із підтвердженим причинним попередником; за повністю спостереженого $X_b=\varnothing$ установлюють $q_t=1$, оскільки непідтверджених записів немає. У дослідній конфігурації прийнято $\mu_A=(m_s+m_v+m_b)/3$ і $c_I=\mu_A\min(q_p,q_t)$ лише за $m_b=1$. За $m_b=0$ значення $c_I$ недоступне, а статичну упевненість формують окремо без поведінкового внеску. Використання мінімуму не дає одному показнику покриття компенсувати інший. Низька упевненість при цьому не зменшує статичного ризику автоматично.

Показники $q_I$ і $c_I$ зберігають у супровідних відомостях $C_I$ інтегрованого представлення. До перевірки спільного калібрування в $F_D$ під час обчислення підсумкової упевненості $c_A$ враховують $c_I$, тоді як $q_I$ залишають діагностичним показником. Вони не замінюють ні $Q_A$, ні інтегрального ризику $\rho_A$.

\subsection{Експериментально перевірена скорочена конфігурація}

Повний метод потребує причинно пов’язаного журналу $X_b$, представлення $Z_b$ і оцінки прояву $Q_A$. Цих даних у виконаних серіях не було. Водночас у дослідному пакеті збережено окрему скорочену конфігурацію, яка узгоджує три статичні оцінки першого методу з урахуванням їх розбіжностей і формату мультимедійного об’єкта. Вона не формує $X_b$ або $Z_b$, не встановлює причинного зв’язку між подіями та статичними координатами й тому не замінює повної конфігурації другого методу. Її призначення полягає у перевірці того, чи може зафіксоване нелінійне узгодження доступних статичних оцінок зменшити кількість хибнопозитивних рішень на нових пристроях.

Нехай $\rho_s$, $\rho_v$ і $\rho_c$ --- відповідно структурна оцінка, оцінка спектральних і піксельних ознак та спільна статична оцінка першого методу, а $e_f\in\{0,1\}^{5}$ --- вектор формату, одна одинична компонента якого позначає один із форматів: PNG, JPEG, WebP, MP4 або WebM. Десятикомпонентний вхід скороченої конфігурації визначено формулою~\eqref{eq:ch3_reduced_alignment_input}:
\begin{equation}
\begin{gathered}
x_2=\bigl[\rho_s,\rho_v,\rho_c,|\rho_s-\rho_v|,|\rho_c-\rho_s|,e_f^{\mathsf T}\bigr]^{\mathsf T},\\
e_f\in\{0,1\}^{5},\qquad \sum_{k=1}^{5}e_{f,k}=1,
\end{gathered}
\label{eq:ch3_reduced_alignment_input}
\end{equation}

де $x_2\in\mathrm R^{10}$ --- вхідний вектор скороченої конфігурації. Ознаки формату не є поведінковими свідченнями; вони лише дають змогу врахувати встановлену в експериментах залежність статичних оцінок від способу подання мультимедійного об’єкта.

Після стандартизації за середніми значеннями $\mu_k$ і масштабами $s_k$, визначеними тільки на навчальній частині, шість прихованих компонентів і вихідну оцінку обчислюють за формулою~\eqref{eq:ch3_reduced_alignment_network}:
\begin{equation}
\begin{gathered}
\widetilde{x}_{2,k}=\frac{x_{2,k}-\mu_k}{s_k},\qquad k=1,\ldots,10,\\
h=\tanh\!\left(W_1^{\mathsf T}\widetilde{x}_2+b_1\right),\qquad
p_2=\sigma\!\left(w_2^{\mathsf T}h+b_2\right),
\end{gathered}
\label{eq:ch3_reduced_alignment_network}
\end{equation}

де $W_1\in\mathrm R^{10\times6}$, $b_1\in\mathrm R^6$, $w_2\in\mathrm R^6$ і $b_2\in\mathrm R$ --- визначені параметри; $h\in\mathrm R^6$ --- внутрішнє представлення; $\sigma(a)=(1+\exp(-a))^{-1}$; $p_2\in[0,1]$ --- оцінка скороченої конфігурації. Разом мережа містить $10\cdot6+6+6+1=73$ визначені параметри.

Збережене правило враховує поріг мережі та окрему двійкову ознаку технічної аномалії контрольованого розкодування. Його подано формулою~\eqref{eq:ch3_reduced_alignment_decision}:
\begin{equation}
\widehat y_2=\mathbf 1\!\left[p_2\geq\tau_2\ \lor\ a_d=1\right],
\label{eq:ch3_reduced_alignment_decision}
\end{equation}

де $\tau_2=0{,}5911823092477697$ --- поріг, визначений на окремій перевірній частині; $a_d\in\{0,1\}$ --- зареєстрована ознака технічної аномалії розкодування; $\widehat y_2$ --- рішення скороченої конфігурації. Значення $a_d$ не ототожнюють із $m_b$, $\rho_b$ або $Q_A$: воно не підтверджує наявності причинно пов’язаного журналу. У зовнішній серії $a_d=0$ для всіх $800$ відеоматеріалів, тому одержані рішення визначалися лише $p_2$.

Параметри визначали за $2\,500$ об’єктами, утвореними із записів пристрою D24 у приміщенні. Ще $2\,500$ об’єктів з іншої сцени того самого пристрою використали лише для вибору порога за найбільшою збалансованою часткою правильних рішень із подальшою перевагою більшої специфічності та більшого порога. Внутрішню перевірку провели на $5\,000$ об’єктах пристрою D27. Використано початкове значення генератора псевдовипадкових чисел $20\,260\,717$, $3\,000$ ітерацій із залученням усієї навчальної частини, швидкість навчання $0{,}01$, коефіцієнт $L_2$-регуляризації $0{,}001$ та зважену за класами бінарну перехресну ентропію. Об’єкти пристроїв D28--D31 не використовували для визначення параметрів або порога.

Машинозчитуваний паспорт із будовою, параметрами, статистиками стандартизації та порогом збережено як \nolinkurl{data/method_2_neural/model.json}; його контрольна сума SHA-256 дорівнює \nolinkurl{32326de5129b8fd1e678060dfa8b42c7451ba99f6479612e04aded066e959156}. Автоматична перевірка підтвердила відсутність перетину адрес джерел і контрольних сум між частиною розроблення та заздалегідь відокремленою зовнішньою серією, а також нульові перетини базових фрагментів між навчальною, перевірною і внутрішньою перевірною частинами. Результати застосування зафіксованої моделі до $2\,000$ похідних від $100$ базових фрагментів і чотирьох нових пристроїв наведено в підрозділі~\ref{sec:ch4_ind_svs}. Ці $2\,000$ об’єктів не трактують як $2\,000$ незалежних джерел.

Межу між формальним складом методу та реально наявними доказами систематизовано в табл.~\ref{tab:ch3_method2_evidence_boundary}. Вона є частиною опису відтворюваності: відсутній складник не можна вважати реалізованим лише тому, що для нього визначено позначення або формулу.

\begin{table}[htbp]
\centering
\small
\caption{Наявність складників, потрібних для перевірки другого методу}
\label{tab:ch3_method2_evidence_boundary}
\begin{tabularx}{\textwidth}{@{}>{\raggedright\arraybackslash}p{0.33\textwidth}>{\raggedright\arraybackslash}p{0.18\textwidth}Y@{}}
\toprule
Складник & Стан & Межа підтвердження \\
\midrule
Статичні оцінки $\rho_s$, $\rho_v$ і $\rho_c$ & Наявні & Відтворено їх обчислення зі збережених ознак; межі першого методу залишаються чинними \\
Причинно пов’язаний журнал $X_b$ із фазами та попередниками & Відсутній & Відбір $P_b$ і повне представлення $Z_b$ експериментально не перевірено \\
Словники, повний каталог правил $\theta_r$ і пороги $\theta_{\rho b}$ & Неповні & Конфігурацію $Z_b=Z_r$ визначено формально, але її повторне виконання за наявним корпусом неможливе \\
Послідовнісне кодування $U_b$ і перехресне зіставлення $A$ & Відсутні & Розширення на основі навченої моделі не входить до підтвердженого результату \\
Скорочена модель узгодження $x_2\mapsto p_2$ & Наявна & Збережено $73$ параметри, поріг, код повторного визначення параметрів і заздалегідь відокремлений протокол \\
Зовнішня серія на D28--D31 & Наявна & Перевірено узгодження статичних оцінок і перенесення між форматами, але не повний контур причинного узгодження ознак подій і поведінки \\
\bottomrule
\end{tabularx}
\end{table}

Отже, емпіричне підтвердження другого методу має два різні рівні. Для повної конфігурації підтверджено формальну визначеність входів, виходів, граничних станів і порядку майбутньої перевірки, але не її кількісні властивості. Для скороченої конфігурації підтверджено відтворювану будову, зафіксовані параметри та зовнішнє оцінювання узгодження статичних оцінок. Позитивний результат другого рівня не переносять на відсутні $X_b$, $Z_b$, $\rho_b$ і $Q_A$.

\subsection{Властивості, обчислювальна складність і перевірка методу}

Для конфігурації з явними правилами визначено чотири властивості, виконання яких необхідне незалежно від числових значень оцінок. По-перше, за незмінних $H$, $J_b$, $C_I$ і версій конфігурації повторне виконання має формувати той самий $X_b$, ту саму маску й той самий результат $Z$ або $Z^{\partial}$. По-друге, кожна подія у $X_b$ повинна задовольняти умову $d(e)=1$; ця властивість забезпечує причинну замкненість відібраної послідовності. По-третє, недоступна група даних не заміщується нульовим вектором і не може зменшити оцінку ризику лише через технічну відсутність спостереження. По-четверте, кожний числовий компонент $Z_b$ має бути простежуваним до правила, унормованої події та початкового запису журналу.

Зазначені властивості не доводять правильності розпізнавання аномального вмісту, але встановлюють формальні умови відтворюваності методу. Додавання до $J_b$ запису, що не пройшов причинно-фазовий допуск, не повинно змінювати $X_b$ і $Z_b$. Натомість додавання допустимої події може змінити результат лише в новому запуску з новим посиланням на журнал. Метод не створює відсутню модальність і не переносить значення з одного запуску до іншого.

Нехай $N_J$ --- кількість записів у початковому журналі, $T$ --- кількість допущених подій, $n_r$ --- кількість правил, які перевіряють для однієї події, а $d_s$, $d_v$ і $d_b$ --- розмірності відповідно представлення структурних ознак, представлення спектральних і піксельних ознак та представлення ознак подій і поведінки. За одноразового перегляду журналу, упорядкування допущених подій та послідовного застосування каталогу верхню межу часової складності визначено формулою~\eqref{eq:ch3_behavior_complexity}:
\begin{equation}
C_M=O\!\left(N_J+T\log T+Tn_r+d_s+d_v+d_b\right),
\label{eq:ch3_behavior_complexity}
\end{equation}

де доданок $N_J$ відповідає відбору записів, $T\log T$ --- упорядкуванню за часовими та причинними відношеннями, $Tn_r$ --- застосуванню правил, а сума розмірностей --- формуванню інтегрованого представлення. Якщо журнал уже впорядкований і має індекси об’єкта та запуску, доданок упорядкування може бути зменшений, однак це не змінює вимог до перевірки причинного попередника. Верхня межа пам’яті становить $O(Td_e+d_s+d_v+d_b)$, де $d_e$ --- розмір числового опису однієї події. Наведені оцінки характеризують порядок зростання витрат, а не виміряну швидкодію відсутньої повної реалізації.

Порядок майбутньої перевірки повної конфігурації має охоплювати не лише підсумкові показники класифікації, а й порушення кожної передумови методу. Мінімальну сукупність перевірних сценаріїв подано в табл.~\ref{tab:ch3_behavior_verification}. Очікувані результати сформульовано до одержання експериментальних даних, тому таблиця є протоколом перевірки, а не повідомленням про вже проведений експеримент.

\begin{table}[htbp]
\centering
\footnotesize
\caption{Перевірні сценарії методу поведінкового узгодження}
\label{tab:ch3_behavior_verification}
\begin{tabularx}{\textwidth}{@{}>{\raggedright\arraybackslash}p{0.25\textwidth}Y>{\raggedright\arraybackslash}p{0.28\textwidth}@{}}
\toprule
Сценарій & Контрольована зміна & Очікуваний результат \\
\midrule
Повторне виконання & Незмінні входи та версії конфігурації & Побайтово однакові $X_b$, маска, $Z_b$ і підсумковий запис \\
Внесення сторонньої події & До журналу додано запис іншого об’єкта, запуску або причинного ланцюга & Запис не входить до $X_b$ і не змінює результат \\
Відсутність обов’язкової фази & Вилучено всі записи однієї обов’язкової фази & $m_b=0$, формується $Z^{\partial}$ із зазначенням неповного покриття \\
Спостережуваний порожній журнал & Усі фази покриті, але допустимих подій немає & $m_b=1$, $X_b=\varnothing$; стан не ототожнюється з недоступністю \\
Порушення причинного порядку & Змінено часові позначки або посилання на попередника & Подію вилучено або сформовано стан відмови згідно з правилами покриття \\
Несумісна версія схеми & Змінено значення версії без відповідного словника перетворення & Автоматичне кодування заборонено; причина несумісності зареєстрована \\
Неперевірене навчуване розширення & Доступні значення $U_b$ або $A$, але немає незалежно перевіреної версії моделі & Значення не впливають на $Z_b$ і підсумкове рішення \\
\bottomrule
\end{tabularx}
\end{table}

Повна емпірична перевірка потребує парних журналів штатного й контрольовано зміненого оброблення тих самих об’єктів, наперед визначених схем подій, каталогу правил, правил причинного зв’язування та незалежних груп для налаштування й оцінювання. Спочатку потрібно перевірити відтворюваність $P_b$ і правильність станів доступності, після цього --- якість $Z_b$, а лише потім --- внесок повної інтеграції у достовірність рішення. Така послідовність не дозволяє приписати другому методу поліпшення, спричинене лише статичними оцінками першого методу.

Алгоритм конфігурації з явними правилами формалізовано, однак його фактичне повторення потребує каталогу правил, унормованих подій, словників, порядку компонентів і порогів. Повного комплекту цих матеріалів у наявному корпусі не зафіксовано. Також відсутні повний набір причинно пов’язаних послідовностей подій, навчений послідовнісний засіб кодування, налаштоване перехресне зіставлення та незалежне порівняння.

Отже, у підрозділі розроблено метод із визначеними метою, об’єктом опрацювання, описами входу й виходу, послідовністю перетворень, умовами допустимості, граничними станами, властивостями та порядком перевірки. Наявні матеріали підтверджують формальну завершеність опису, але не кількісну перевагу повної конфігурації. Інформаційна технологія може організувати застосування конфігурації з явними правилами та зберегти відомості про її стан, проте не замінює окремої експериментальної перевірки другого методу.

\section{Інформаційна технологія виявлення ознак аномального вмісту в мультимедійних даних вебсередовища}
\label{sec:information_technology}

Окреме виконання двох методів ще не визначає, коли об’єкт можна передати вебсистемі, як діяти за неповного результату та як відтворити підстави рішення після зміни параметрів. Технологічне завдання полягає в керуванні станами недовіреного носія, а не лише в послідовному виконанні обчислень. До завершення аналізу об’єкт потрібно зберігати в ізоляції, після аналізу --- відокремити аналітичний стан від припису за правилом реагування, а після визначення припису --- зареєструвати відомості, достатні для його перевірки. Зовнішню дію виконують лише після успішної реєстрації.

Інформаційна технологія використовує одну модель і два методи, але не змінює їхнього наукового змісту. Модель задає структуру аналітичного стану $S_A$, перший метод формує статичний опис, а другий --- інтегроване представлення. Технологія керує ізоляцією, порядком виконання, вибором припису, реконструкцією та реєстрацією. За такого розмежування правило реагування можна змінювати без зміни висновку моделі, а перевірку --- повторювати без заміни початкового об’єкта реконструйованим.

Наукова відмінність технології полягає не в самому числі етапів, а в незмінних вимогах до процесу: збереженні початкового $X_{mm}$; станах із визначеним складом та явною маскою доступності; відокремленні аналітичного стану $S_A$ від технологічного припису $D_A$; поданні реконструкції як нового циклу з власним ідентифікатором, а не як «очищення» того самого носія; реєстраційному записі, достатньому для повторної перевірки за умови гарантованого збереження початкових даних і журналу. Склад п’яти етапів установлено в підрозділі~\ref{sec:process_organization}; нижче конкретизовано правило реагування, реєстрацію та реконструкцію. Зіставлення з аналогами наведено в табл.~\ref{tab:ch3_novelty_comparison}.

П’ятиетапну послідовність і забезпечувальні компоненти показано на рис.~\ref{fig:information_technology}. Засіб приймання, ізольоване сховище, контрольоване середовище та сховище записів не утворюють додаткових етапів і не змінюють логіки переходів.

\begin{figure}[htbp]
\centering
\begin{tikzpicture}[x=1cm,y=1cm]
\node[dissertation block,text width=3.65cm] (s1) at (-4.70,1.75)
  {1. Попереднє оброблення\\$S_0\rightarrow S_1$\\$X_{mm}\rightarrow X_c$};
\node[dissertation block,text width=3.65cm] (s2) at (0,1.75)
  {2. Формування ознак\\$S_1\rightarrow S_2$\\$H,J_b$};
\node[dissertation block,text width=3.65cm] (s3) at (4.70,1.75)
  {3. Відбір та інтеграція\\$S_2\rightarrow S_3$\\$X_b,Z$};
\node[dissertation block,text width=4.15cm] (s4) at (2.40,-1.40)
  {4. Прийняття рішення та визначення припису\\$S_3\rightarrow S_4$\\$S_4=\langle S_A,D_A\rangle$};
\node[dissertation block,text width=4.15cm] (s5) at (-2.40,-1.40)
  {5. Реєстрація результатів\\$S_4\rightarrow S_5$\\$S_5=F_R(S_4,C_R)$};

\draw[dissertation arrow] (s1.east) -- (s2.west);
\draw[dissertation arrow] (s2.east) -- (s3.west);
\draw[dissertation arrow] (s3.south) -- (s4.north);
\draw[dissertation arrow] (s4.west) -- (s5.east);
\end{tikzpicture}
\caption{П’ять етапів інформаційної технології виявлення ознак аномального вмісту в мультимедійних даних вебсередовища}
\label{fig:information_technology}
\end{figure}


Інформаційна технологія $T_A$ охоплює впорядковані етапи $T_1,\ldots,T_5$. Пропуск будь-якого з них змінює зміст результату: стан після прийняття рішення ще не є завершеним технологічним результатом, доки не збережено відомостей для його відтворення. Засіб приймання, ізольоване сховище, контрольоване середовище, засіб реконструкції та сховище записів є забезпечувальними компонентами. Для кожного переходу між ними визначено вхід, склад виходу, контрольну суму, версію та стан завершення.

Аналітичний результат не визначає припису без урахування умов застосування: однаковий ризик може бути прийнятним для внутрішнього архіву й неприйнятним для публічного каналу. Щоб ця відмінність не змінювала модельного стану заднім числом, правило реагування відокремлено від $F_D$. Залежність технологічного припису від аналітичного стану, середовища застосування й чинної версії правила визначено формулою~\eqref{eq:ch3_policy_action}:
\begin{equation}
D_A=F_A(S_A,C_p,\theta_p),
\label{eq:ch3_policy_action}
\end{equation}

де $D_A$ --- технологічний припис; $F_A$ --- функція вибору припису за правилом реагування; $C_p$ --- незмінний опис середовища застосування; $\theta_p$ --- параметри та версія цього правила; $S_A$ --- аналітичний стан.

Допустимими приписами є допуск, передавання на експертну перевірку, реконструкція та блокування. За високого ризику й низької упевненості правило не повинно автоматично визначати допуск. Однаковий $S_A$ може спричинити різні приписи в різних середовищах, але припис не змінює модельного висновку. Формула~\eqref{eq:ch3_policy_action} дає змогу змінювати правила реагування без повторного тлумачення вже сформованого аналітичного стану.

Після застосування правила реагування утворюється стан четвертого етапу, але повний цикл завершується лише після успішної реєстрації. Склад цього проміжного стану, у якому аналітичний результат і припис залишаються окремими складниками, визначено формулою~\eqref{eq:ch3_output_state}:
\begin{equation}
S_4=\langle S_A,D_A\rangle,
\label{eq:ch3_output_state}
\end{equation}

де $S_4$ --- стан після етапу прийняття рішення; $S_A$ --- аналітичний результат; $D_A$ --- технологічний припис.

За такого розділення окремо простежують наслідки зміни ознак, параметрів моделі та правила реагування.

Збереження лише $S_4$ недостатньо для пояснення зміни результату після оновлення засобу розкодування, схеми ознак або правила реагування. Тому на п’ятому етапі до цього стану приєднують відомості для відтворення й утворюють зареєстрований технологічний результат за формулою~\eqref{eq:ch3_audit_record}:
\begin{equation}
S_5=F_R(S_4,C_R),
\label{eq:ch3_audit_record}
\end{equation}

де $F_R$ --- функція реєстрації; $S_5$ --- зареєстрований технологічний результат; $S_4$ --- стан після прийняття рішення; $C_R$ --- відомості, потрібні для відтворення.

Відомості $C_R$ містять ідентифікатор і контрольну суму $X_{mm}$, версії засобів, схем ознак і правила реагування, опис $C_p$, ідентифікатор запуску, маску $m_A$, координати $L_s$ і $L_v$, часові межі та причини неповноти. Представлення $X_b$ зберігають як незмінну копію або як адресоване за вмістом посилання на збережений журнал; у другому разі до $C_R$ включають контрольну суму, строк зберігання та гарантію доступності. Без цих умов запис є достатнім для повторної перевірки лише за фактичної доступності журналу.

Реєстрація завершує формування технологічного результату. Зовнішній припис $D_A$ виконують лише після цілісного створення $S_5$; допуск до публікації до цього моменту заборонено. Якщо запис не створено, припис не виконують, об’єкт залишають в ізоляції, а система може призначити повторне опрацювання, експертний розгляд або блокування. Виконання здійснюють ідемпотентно за ідентифікатором запуску, а його підтверджений результат додають до реєстраційних відомостей. Перерваний запуск не перезаписує завершеного запису; повторне виконання одержує новий ідентифікатор.

Контрольована реконструкція потрібна тоді, коли чинне правило реагування допускає збереження дозволеного піксельного вмісту, але початкова структура контейнера не вважається надійною. Вона створює новий мультимедійний об’єкт без невідомих невіднесених ділянок і заборонених службових полів; це не очищення початкового носія, а формування іншого об’єкта. Відповідне перетворення встановлено формулою~\eqref{eq:ch3_reconstruction_candidate}:
\begin{equation}
X_{mm,j}=E_R(X_{v,i},\theta_R),\quad j\neq i,
\label{eq:ch3_reconstruction_candidate}
\end{equation}

де $E_R$ --- функція реконструкції; $X_{v,i}$ --- дозволене піксельне або кадрове подання початкової перевірки; $\theta_R$ --- опис цільового формату; $X_{mm,j}$ --- новий мультимедійний об’єкт.

Реконструкція у формулі~\eqref{eq:ch3_reconstruction_candidate} не замінює $X_{mm,i}$ і не надає новому об’єкту наперед установленого статусу безпечності. Зв’язок із джерелом зберігають у супровідних відомостях, але новий об’єкт одержує власний ідентифікатор. Для відео перетворення охоплює лише кадрове подання: звукові доріжки, субтитри та службові поля доріжок до нового об’єкта автоматично не переносять, а їх склад визначають описом цільового формату $\theta_R$ і чинним правилом реагування.

Оскільки реконструкція змінює байтову послідовність, попередній клас або ризик не можна успадковувати як властивість нового об’єкта. Реконструйований об’єкт повертається на початок і проходить повний технологічний цикл за формулою~\eqref{eq:ch3_reconstruction_cycle}:
\begin{equation}
S_{5,j}=F_T(X_{mm,j},C_{0,j},\theta_T),
\label{eq:ch3_reconstruction_cycle}
\end{equation}

де $F_T$ --- функція виконання п’яти етапів технології; $C_{0,j}$ --- відомості про надходження реконструйованого об’єкта до нового циклу; $\theta_T$ --- зафіксована конфігурація; $S_{5,j}$ --- зареєстрований результат нового п’ятиетапного циклу.

Новий цикл у формулі~\eqref{eq:ch3_reconstruction_cycle} не успадковує класу, ризику або припису попередньої перевірки. Допуск можливий лише за результатом нового аналітичного стану й успішної реєстрації. Реконструкція при цьому не утворює шостого етапу технології: вона не продовжує поточної перевірки, а породжує новий об’єкт і, відповідно, новий п’ятиетапний цикл із власним ідентифікатором.

Ізоляцію об’єкта до зареєстрованого рішення та повернення лише реконструйованого об’єкта на початок нового циклу показано на рис.~\ref{fig:zero_trust}. Тому зворотний зв’язок на схемі не означає повторного використання попереднього рішення.

\begin{figure}[htbp]
\centering
\begin{tikzpicture}[x=1cm,y=1cm]
\node[dissertation block,text width=4.20cm] (object) at (0,4.85)
  {Початковий об’єкт\\$X_{mm}$};
\node[dissertation block,text width=4.20cm] (isolation) at (0,3.25)
  {Ізоляція об’єкта};
\node[dissertation block,text width=4.20cm] (cycle) at (0,1.65)
  {Формування аналітичного стану\\$S_A$};
\node[dissertation block,text width=4.20cm] (registered) at (0,0.05)
  {Політика реагування\\$D_A=F_A(S_A,\theta_p)$};
\node[dissertation block,text width=4.20cm] (policy) at (0,-1.55)
  {Зареєстрований результат\\$S_5=F_R(S_4,C_R)$};

\node[dissertation block,text width=2.45cm] (allow) at (-5.10,-4.00)
  {Допуск};
\node[dissertation block,text width=2.85cm] (review) at (-1.75,-4.00)
  {Експертна перевірка};
\node[dissertation block,text width=2.45cm] (block) at (1.75,-4.00)
  {Блокування};
\node[dissertation block,text width=3.35cm] (reconstruct) at (5.15,-4.00)
  {Реконструкція кандидата\\$X_{mm,j}$};

\draw[dissertation arrow] (object.south) -- (isolation.north);
\draw[dissertation arrow] (isolation.south) -- (cycle.north);
\draw[dissertation arrow] (cycle.south) -- (registered.north);
\draw[dissertation arrow] (registered.south) -- (policy.north);

\coordinate (decisionbus) at (0,-2.80);
\draw[dissertation arrow] (policy.south) -- (decisionbus);
\draw[thick] (-5.10,-2.80) -- (5.15,-2.80);
\draw[dissertation arrow] (-5.10,-2.80) -- (allow.north);
\draw[dissertation arrow] (-1.75,-2.80) -- (review.north);
\draw[dissertation arrow] (1.75,-2.80) -- (block.north);
\draw[dissertation arrow] (5.15,-2.80) -- (reconstruct.north);

\coordinate (loopright) at (7.15,3.25);
\draw[thick] (reconstruct.east) -- (7.15,-4.00) -- (loopright);
\node[font=\small,anchor=east] at (6.95,-0.85) {новий цикл};
\draw[dissertation arrow] (loopright) -- (isolation.east);
\end{tikzpicture}
\caption{Контур нульової довіри для інформаційної технології}
\label{fig:zero_trust}
\end{figure}


Синхронний, асинхронний і фоновий режими не змінюють п’яти етапів. Перевищення обмежень часу чи пам’яті, помилка розкодування та відмова сховища є технічними станами, а не свідченнями нульового ризику. Технічна невдача не є підставою для автоматичної публікації об’єкта. У локальному прототипі, перевіреному в розділі~\ref{chap:experiments}, помилка розкодування й контрольовано внесена відмова реєстрації завершували цикл без створення $S_5$ та без виконання зовнішнього припису; вебобгортка зберігає такий вхід в ізольованому сховищі.

Алгоритмічна відтворюваність технології потребує незмінних входів, параметрів із визначеними версіями, унормованих схем подій, каталогу правил, порогів, ваг і повних реєстраційних записів. Формальна наявність цих вимог, однак, не доводить їх реалізації в розподіленому середовищі. Локальна перевірка підтвердила цілісність переходів $S_1$--$S_5$, атомарний запис без перезапису, повторний запуск, відновлення після двох класів технічних відмов і завершеність одного чотирипроцесного запуску. Не перевірено реконструкцію, причинно пов’язаний журнал, розподілену чергу, мережеву взаємодію, багатовузлове сховище, фактичне виконання зовнішніх приписів і тривале навантаження. Отже, реалізовано й локально перевірено прототипну послідовність виконання інформаційної технології, але її промислову готовність не заявлено.

\section*{Висновки до розділу 3}
\phantomsection
\addcontentsline{toc}{section}{Висновки до розділу 3}

У розділі розв’язано завдання переходу від модельного опису до послідовності операцій зі збереженням походження кожного свідчення. Незмінний вхід відокремлено від похідних представлень, статичний опис --- від початкового журналу й причинно пов’язаних подій, а аналітичний стан --- від технологічного припису. Перший метод формує статичний опис $H$, другий --- інтегроване представлення $Z$, після чого інформаційна технологія організовує їх п’ятиетапне застосування та реєстрацію. Емпірично оцінено окремі статичні аналітичні складники, а локальний прототип перевірено щодо переходів, реєстрації, повтору, відновлення й двох класів технічних відмов. Повну конфігурацію поведінкового узгодження та розподілений промисловий процес не перевірено, тому промислову готовність технології не заявлено.

Удосконалено метод аналізу спектральних, піксельних і структурних ознак шляхом відокремленого формування двох складових ознак зі збереженням байтових і просторово-часових координат. Визначено вхід, параметри, одинадцять послідовних кроків алгоритму, умови завершення, структурні перевірки, детерміновані локальні й частотні ознаки, кадровий відбір та агрегування. Результат $H$ є проміжним статичним описом і не замінює інтегрованого оцінювання або технологічного рішення. За протоколом із груповим поділом у підрозділі~\ref{sec:ch4_ind_svs} перевірено скорочені числові проєкції $Z_s$, $Z_v$ і $Z_c$, але не локалізації, покадрове зважування або повний склад $H$. Результати розділу~\ref{chap:experiments} водночас засвідчили, що ізольована складова спектральних і піксельних ознак на реальних відеозаписах може втрачати роздільну здатність. Отже, призначення методу полягає у спільному розгляді структурних свідчень і результатів спектрального та піксельного аналізу разом із координатами та маскою доступності, а не в автономному використанні спектрального аналізу.

Формалізовано метод поведінкового узгодження ознак різного походження, за яким подія має належати тому самому запуску, дозволеному причинному ланцюгу та сумісній фазі. Для конфігурації з явними правилами визначено каталог відповідностей із установленою версією. Кодування послідовностей і перехресне зіставлення на основі навченої моделі відокремлено як неперевірене розширення. Оскільки повний каталог правил і повні послідовності подій у наявному корпусі не збережено, експериментальну перевагу та кількісні властивості повної конфігурації не встановлено. Скорочена складова, адаптована до формату й навчена за експериментальними даними, яку розглянуто в розділі~\ref{chap:experiments}, не замінює такої перевірки. Отже, назва «метод» позначає завершену формальну процедуру з визначеними входом, перетвореннями, виходом і граничними станами, а не твердження про доведену емпіричну перевагу її повної реалізації.

Набула подальшого розвитку інформаційна технологія, незмінними вимогами якої є збереження входу, визначена структура станів, відокремлення аналітичного стану $S_A$ від припису $D_A$ у стані $S_4=\langle S_A,D_A\rangle$, створення реєстраційного запису $S_5$ до зовнішнього виконання припису й подання реконструкції як нового циклу з окремим ідентифікатором. Локальна перевірка у розділі~\ref{chap:experiments} підтвердила порядок $S_1$--$S_5$, незмінність входу, атомарну реєстрацію, повтор, відновлення та безпечне завершення за помилки розкодування й відмови реєстрації. Цей результат не підтверджує повної поведінкової конфігурації, реконструкції, виконання зовнішніх приписів або роботи технології в розподіленому промисловому середовищі.
